\documentclass[12pt,a4paper]{ctexart}
\usepackage{geometry}
\usepackage{graphicx}
\usepackage{amsmath,amssymb}
\usepackage{booktabs}
\usepackage{array}
\usepackage{enumitem}
\usepackage{xcolor}
\usepackage{listings}
\usepackage{hyperref}
\usepackage{fancyhdr}
\usepackage{tcolorbox}
\tcbuselibrary{skins,breakable}

\geometry{left=2cm,right=2cm,top=2cm,bottom=2cm}
\hypersetup{colorlinks=true,linkcolor=blue}
\lstset{basicstyle=\ttfamily\small,breaklines=true,frame=single,backgroundcolor=\color{gray!10}}

% 设置目录深度（显示到subsection级别）
\setcounter{tocdepth}{2}
\setcounter{secnumdepth}{3}

\newtcolorbox{keybox}{colback=blue!5,colframe=blue!75!black,title=重点,fonttitle=\bfseries,breakable}
\newtcolorbox{cmdbox}{colback=green!5,colframe=green!75!black,title=命令,fonttitle=\bfseries,breakable}
\newtcolorbox{exambox}[1][]{colback=yellow!5,colframe=orange!75!black,title=#1,fonttitle=\bfseries,breakable}

\title{\Huge\textbf{计算机网络实验复习}}
\author{考试知识点整理}
\date{\today}

\begin{document}
\maketitle

% 生成详细目录
\tableofcontents
\newpage

%========================================
\section{网络设备基础}
%========================================

\subsection{集线器(HUB)与交换机(Switch)}

\begin{keybox}
\textbf{集线器的特点}（常考易错）：
\begin{itemize}[itemsep=0pt]
    \item 集线器工作在\textbf{物理层}，不是基于MAC地址转发
    \item 连接到集线器的节点发送数据时，执行\textbf{CSMA/CD}介质访问控制
    \item 通过在网络链路中串接集线器可以\textbf{监听}该链路中的数据包
    \item 连接到同一个集线器的所有节点共享\textbf{同一个冲突域}
\end{itemize}
\end{keybox}

\begin{table}[htbp]
\centering
\begin{tabular}{lcc}
\toprule
\textbf{特性} & \textbf{集线器} & \textbf{交换机} \\
\midrule
工作层次 & 物理层 & 数据链路层 \\
转发依据 & 无（广播） & MAC地址表 \\
冲突域 & 所有端口共享 & 每端口独立 \\
通信方式 & 半双工 & 全双工 \\
\bottomrule
\end{tabular}
\caption{集线器与交换机对比}
\end{table}

\subsection{二层交换机与三层交换机}

\begin{itemize}[leftmargin=2em,itemsep=0pt]
    \item \textbf{二层交换机}：根据\textbf{MAC地址}转发数据帧
    \item \textbf{三层交换机}：二层交换+三层路由，可实现VLAN间路由
    \item \textbf{路由器}：根据\textbf{路由表}和\textbf{目的IP地址}转发IP分组
\end{itemize}

\subsection{交换机管理方式}

\begin{itemize}[leftmargin=2em,itemsep=0pt]
    \item \textbf{Console口}：通过串口线本地配置
    \item \textbf{Telnet}：远程命令行配置
    \item \textbf{Web管理}：通过浏览器访问交换机的\textbf{IP地址}进行配置
\end{itemize}

%========================================
\section{VLAN虚拟局域网}
%========================================

\subsection{VLAN的作用}

\begin{keybox}
\textbf{VLAN的好处}：
\begin{itemize}[itemsep=0pt]
    \item 广播可以得到控制（减小广播域）
    \item 网络安全可以得到保护
    \item 实现虚拟工作组
    \item 灵活的网络管理
\end{itemize}
\textbf{注意}：增加VLAN\textbf{不能}使交换机不需要再配置！
\end{keybox}

\subsection{VLAN划分方法}

\begin{itemize}[leftmargin=2em,itemsep=0pt]
    \item \textbf{基于端口划分}（实验课采用的方法）
    \item 基于MAC地址划分（主机位置变动时无需重新分配）
    \item 基于IP地址划分
    \item 基于协议划分
\end{itemize}

\subsection{VLAN配置命令}

\begin{cmdbox}
\textbf{创建VLAN}：
\begin{lstlisting}
Switch(config)# vlan 2
! 2: VLAN编号(1-4094)
Switch(config-vlan)# name vlan2
! vlan2: VLAN名称(可选)
Switch(config-vlan)# exit
\end{lstlisting}

\textbf{将端口划入VLAN}：
\begin{lstlisting}
Switch(config)# interface fastEthernet 0/1
Switch(config-if)# switchport access vlan 2
! access: 接入模式(连接终端设备)
! vlan 2: 将该端口划入VLAN 2
Switch(config-if)# exit
\end{lstlisting}

\textbf{配置Trunk端口}（交换机互连）：
\begin{lstlisting}
Switch(config)# interface fastEthernet 0/24
Switch(config-if)# switchport mode trunk
! trunk: 中继模式(连接交换机,允许多个VLAN通过)
Switch(config-if)# exit
\end{lstlisting}
\end{cmdbox}

\subsection{VLAN间通信}

\begin{keybox}
不同VLAN之间通信需要\textbf{三层设备}（路由器或三层交换机）。

\textbf{三层交换机配置VLAN间路由}：
\begin{lstlisting}
! 创建VLAN接口并配置IP
Switch(config)# interface vlan 2
! vlan 2: 虚拟VLAN接口
Switch(config-if)# ip address 10.1.1.254 255.0.0.0
! 10.1.1.254: IP地址(作为该VLAN的网关)
! 255.0.0.0: 子网掩码
Switch(config-if)# exit

Switch(config)# interface vlan 3
Switch(config-if)# ip address 10.2.1.254 255.0.0.0
Switch(config-if)# exit

! 启用路由功能
Switch(config)# ip routing
! 启用三层交换机的路由功能,实现VLAN间通信
\end{lstlisting}
\end{keybox}

%========================================
\section{IP地址与子网划分}
%========================================

\subsection{子网掩码计算}

\begin{exambox}[常考题型]
\textbf{题目}：一个主机的IP地址是11.0.11.3/24，该主机所属的网络地址和反掩码为？

\textbf{解答}：
\begin{itemize}[itemsep=0pt]
    \item /24表示子网掩码为255.255.255.0
    \item 网络地址 = 11.0.11.3 AND 255.255.255.0 = \textbf{11.0.11.0}
    \item 反掩码（通配符掩码）= 255.255.255.255 - 255.255.255.0 = \textbf{0.0.0.255}
\end{itemize}
\end{exambox}

\begin{exambox}[常考题型]
\textbf{题目}：一个局域网中某台主机的IP地址为176.68.160.12，使用22位作为网络地址，求子网掩码。

\textbf{解答}：
\begin{itemize}[itemsep=0pt]
    \item 22位网络地址 = 11111111.11111111.11111100.00000000
    \item 子网掩码 = \textbf{255.255.252.0}
\end{itemize}
\end{exambox}

\subsection{默认网关}

\begin{keybox}
默认网关地址是\textbf{通往其他网络传输路径上的第一跳路由器地址}。

当主机要访问不在同一网段的目标时，数据包会发送给默认网关。
\end{keybox}

%========================================
\section{路由器配置}
%========================================

\subsection{路由类型}

\begin{itemize}[leftmargin=2em,itemsep=0pt]
    \item \textbf{直连路由}：路由器接口直接连接的网络，自动生成
    \item \textbf{静态路由}：必须由网络管理员\textbf{手动配置}
    \item \textbf{动态路由}：通过路由协议自动学习
\end{itemize}

\subsection{静态路由配置}

\begin{cmdbox}
\textbf{配置静态路由}：
\begin{lstlisting}
! 格式：ip route 目的网络 子网掩码 下一跳地址
Router(config)# ip route 目的网络 子网掩码 下一跳地址
! 目的网络: 要到达的目标网络地址
! 子网掩码: 目标网络的子网掩码
! 下一跳地址: 转发数据包的下一个路由器接口IP地址
\end{lstlisting}

\textbf{配置默认路由}：
\begin{lstlisting}
Router(config)# ip route 0.0.0.0 0.0.0.0 下一跳地址
! 0.0.0.0 0.0.0.0: 匹配所有网络(默认路由)
! 下一跳地址: 当路由表中没有匹配项时,转发到此地址
\end{lstlisting}

\textbf{示例}：
\begin{lstlisting}
Router(config)# ip route 10.10.10.0 255.255.255.0 192.168.11.1
! 到达10.10.10.0/24网络,下一跳为192.168.11.1
Router(config)# ip route 0.0.0.0 0.0.0.0 192.168.10.1
! 默认路由,所有未匹配的数据包转发到192.168.10.1
\end{lstlisting}
\end{cmdbox}

\subsection{路由器接口配置}

\begin{cmdbox}
\textbf{配置以太网接口}：
\begin{lstlisting}
Router(config)# interface fastEthernet 0/0
! interface: 进入接口配置模式
! fastEthernet: 快速以太网接口类型(100Mbps)
! 0/0: 插槽号/端口号(格式: 插槽/端口)
Router(config-if)# ip address 192.168.1.1 255.255.255.0
! ip address: 配置接口IP地址命令
! 192.168.1.1: 接口IP地址
! 255.255.255.0: 子网掩码(/24)
Router(config-if)# no shutdown
! no shutdown: 启用接口(接口默认是关闭状态)
Router(config-if)# exit
! exit: 退出接口配置模式
\end{lstlisting}

\textbf{配置串口}：
\begin{lstlisting}
Router(config)# interface serial 0/0
! serial: 串行接口类型(用于WAN连接)
! 0/0: 插槽号/端口号
Router(config-if)# ip address 192.168.2.1 255.255.255.0
Router(config-if)# clock rate 64000
! clock rate: 配置时钟速率(单位: bps)
! 64000: 64kbps的时钟速率
! 注意: 仅DCE端需要配置,DTE端不需要
Router(config-if)# no shutdown
Router(config-if)# exit
\end{lstlisting}
\end{cmdbox}

\subsection{查看路由表}

\begin{cmdbox}
\begin{lstlisting}
Router# show ip route
\end{lstlisting}
\end{cmdbox}

%========================================
\section{访问控制列表(ACL)}
%========================================

\subsection{ACL类型}

\begin{itemize}[leftmargin=2em,itemsep=0pt]
    \item \textbf{标准ACL}：编号1-99，只能根据\textbf{源地址}过滤
    \item \textbf{扩展ACL}：编号100-199，可根据源地址、目的地址、协议、端口过滤
\end{itemize}

\subsection{ACL配置命令}

\begin{cmdbox}
\textbf{标准ACL}：
\begin{lstlisting}
Router(config)# access-list 1 permit 192.168.1.1
! access-list: 创建访问控制列表
! 1: ACL编号(1-99为标准ACL)
! permit: 允许(deny为拒绝)
! 192.168.1.1: 源IP地址

Router(config)# access-list 1 permit 192.168.1.0 0.0.0.255
! 192.168.1.0: 源网络地址
! 0.0.0.255: 反掩码(wildcard mask),0表示必须匹配,1表示任意
! 此规则匹配192.168.1.0-192.168.1.255

Router(config)# access-list 1 deny any
! any: 匹配所有地址(等同于0.0.0.0 255.255.255.255)
! 注意: ACL末尾有隐含的deny any规则
\end{lstlisting}

\textbf{扩展ACL}（限制Web服务访问）：
\begin{lstlisting}
Router(config)# access-list 101 deny tcp 192.168.1.0 0.0.0.255 any eq 80
! 101: ACL编号(100-199为扩展ACL)
! deny: 拒绝
! tcp: 协议类型(可选: tcp/udp/icmp/ip等)
! 192.168.1.0 0.0.0.255: 源地址和反掩码
! any: 任意目的地址
! eq 80: 等于80端口(HTTP服务)
! eq: 等于, gt: 大于, lt: 小于, neq: 不等于

Router(config)# access-list 101 permit ip any any
! ip: 匹配所有IP协议
! any any: 任意源地址到任意目的地址
\end{lstlisting}

\textbf{应用ACL到接口}：
\begin{lstlisting}
Router(config)# interface fastEthernet 0/0
Router(config-if)# ip access-group 1 in
! ip access-group: 将ACL应用到接口
! 1: ACL编号
! in: 对进入接口的数据流进行过滤
Router(config-if)# ip access-group 1 out
! out: 对离开接口的数据流进行过滤
\end{lstlisting}
\end{cmdbox}

\begin{keybox}
\textbf{ACL规则}：
\begin{itemize}[itemsep=0pt]
    \item 如果分组不匹配当前规则，则检查下一条语句
    \item 如果所有规则都不匹配，则\textbf{隐含拒绝}（deny any）
    \item 标准ACL应用在\textbf{靠近目的地}的位置
    \item 扩展ACL应用在\textbf{靠近源地址}的位置
\end{itemize}
\end{keybox}

%========================================
\section{NAT网络地址转换}
%========================================

\subsection{NAT工作原理}

\begin{keybox}
NAT将内部私有IP地址转换为外部公网IP地址，使内网主机能访问Internet。

\textbf{地址术语}：
\begin{itemize}[itemsep=0pt]
    \item \textbf{内部本地地址}：内网主机的私有IP（如192.168.1.3）
    \item \textbf{内部全局地址}：转换后的公网IP（如59.67.148.3）
\end{itemize}
\end{keybox}

\subsection{NAT转换过程}

\begin{exambox}[NAT转换示例]
内网主机192.168.1.3:3105访问外网服务器202.113.64.2:8080

\textbf{出站数据包}（内→外）：
\begin{itemize}[itemsep=0pt]
    \item 原始：源IP=192.168.1.3:3105，目的IP=202.113.64.2:8080
    \item 转换后：源IP=59.67.148.3:5234，目的IP=202.113.64.2:8080
\end{itemize}

\textbf{入站数据包}（外→内）：
\begin{itemize}[itemsep=0pt]
    \item 原始：源IP=202.113.64.2:8080，目的IP=59.67.148.3:5234
    \item 转换后：源IP=202.113.64.2:8080，目的IP=192.168.1.3:3105
\end{itemize}
\end{exambox}

\subsection{NAT配置命令}

\begin{cmdbox}
\begin{lstlisting}
Router(config)# access-list 11 permit 192.168.1.0 0.0.0.255
! access-list 11: 创建ACL编号11
! permit: 允许转换
! 192.168.1.0 0.0.0.255: 允许转换的内网地址范围

Router(config)# ip nat pool abc 59.67.148.3 59.67.148.3 netmask 255.255.255.0
! ip nat pool: 定义NAT地址池
! abc: 地址池名称(自定义)
! 59.67.148.3 59.67.148.3: 地址池起始IP和结束IP(可以是范围)
! netmask 255.255.255.0: 地址池的子网掩码

Router(config)# ip nat inside source list 11 pool abc overload
! ip nat inside source: 配置内部源地址转换
! list 11: 使用ACL 11定义的地址范围
! pool abc: 使用名为abc的地址池
! overload: 启用PAT(端口地址转换),允许多个内网IP共享一个公网IP

Router(config)# interface fastEthernet 0/0
Router(config-if)# ip nat inside
! ip nat inside: 定义为内部接口(连接内网)

Router(config)# interface serial 0/0
Router(config-if)# ip nat outside
! ip nat outside: 定义为外部接口(连接外网/Internet)
\end{lstlisting}
\end{cmdbox}

%========================================
\section{生成树协议(STP)}
%========================================

\subsection{STP作用}

防止交换机之间形成环路，避免广播风暴。

\subsection{STP选举规则}

\begin{keybox}
\textbf{根交换机选举}：选择\textbf{最小的交换机优先级}（优先级相同则选最小MAC地址）

\textbf{根端口选举}：非根交换机上选择到根交换机\textbf{路径开销最小}的端口（开销相同则选最小端口编号）
\end{keybox}

%========================================
\section{Windows网络命令}
%========================================

\subsection{常用命令}

\begin{cmdbox}
\textbf{查看IP配置}：
\begin{lstlisting}
ipconfig              ! 查看IP地址
ipconfig /all         ! 查看详细信息（包括MAC地址）
\end{lstlisting}

\textbf{测试连通性}：
\begin{lstlisting}
ping 目标IP           ! 测试连通性
ping -t 目标IP        ! 持续测试
ping -n 4 目标IP      ! 指定发送4个包
\end{lstlisting}

\textbf{DNS相关}：
\begin{lstlisting}
nslookup 域名         ! DNS解析
ipconfig /flushdns    ! 清空DNS缓存
\end{lstlisting}

\textbf{ARP相关}：
\begin{lstlisting}
arp -a                ! 查看ARP缓存
arp -d                ! 清空ARP缓存
arp -s IP MAC         ! 手动添加静态ARP条目
\end{lstlisting}

\textbf{路由相关}：
\begin{lstlisting}
route print           ! 查看路由表
tracert 目标IP        ! 路由跟踪
\end{lstlisting}

\textbf{网络统计}：
\begin{lstlisting}
netstat -an           ! 查看网络连接状态
\end{lstlisting}
\end{cmdbox}

\subsection{Ping命令详解}

\begin{keybox}
\textbf{ping命令使用ICMP协议}，类型为Echo Request/Echo Reply。

\textbf{测试顺序}（排查网络故障）：
\begin{enumerate}[itemsep=0pt]
    \item ping 127.0.0.1 — 测试本机TCP/IP协议栈
    \item ping 本机IP — 测试本机网卡
    \item ping 默认网关 — 测试本网段连通性
    \item ping 远程主机 — 测试跨网段连通性
\end{enumerate}
\end{keybox}

%========================================
\section{常用协议与端口}
%========================================

\subsection{协议功能}

\begin{table}[htbp]
\centering
\begin{tabular}{lll}
\toprule
\textbf{协议} & \textbf{功能} & \textbf{端口/协议} \\
\midrule
DNS & 域名解析为IP地址 & UDP 53 \\
ARP & IP地址解析为MAC地址 & — \\
RARP & MAC地址解析为IP地址 & — \\
ICMP & 网络控制报文（ping使用） & — \\
FTP & 文件传输 & TCP 20(数据)/21(控制) \\
HTTP & Web服务 & TCP 80 \\
HTTPS & 安全Web服务 & TCP 443 \\
Telnet & 远程登录 & TCP 23 \\
SMTP & 发送邮件 & TCP 25 \\
\bottomrule
\end{tabular}
\caption{常用协议与端口}
\end{table}

\subsection{FTP命令}

\begin{cmdbox}
\begin{lstlisting}
ftp 服务器IP          ! 连接FTP服务器
ls                    ! 列出文件
get 文件名            ! 下载文件
put 文件名            ! 上传文件
quit                  ! 退出
\end{lstlisting}

\textbf{匿名登录}：用户名输入\texttt{anonymous}
\end{cmdbox}

%========================================
\section{TCP三次握手}
%========================================

\begin{keybox}
\textbf{TCP建立连接的三次握手}：

\begin{center}
\begin{tabular}{|c|c|c|c|c|}
\hline
\textbf{报文} & \textbf{Sequence Number} & \textbf{Ack Number} & \textbf{SYN} & \textbf{ACK} \\
\hline
第一次握手 & x（随机） & 0 & 1 & 0 \\
\hline
第二次握手 & y（随机） & x+1 & 1 & 1 \\
\hline
第三次握手 & x+1 & y+1 & 0 & 1 \\
\hline
\end{tabular}
\end{center}

\textbf{请求连接报文标志}：SYN=1
\end{keybox}

%========================================
\section{双绞线制作}
%========================================

\subsection{线序标准}

\begin{itemize}[leftmargin=2em,itemsep=0pt]
    \item \textbf{T568A}：白绿、绿、白橙、蓝、白蓝、橙、白棕、棕
    \item \textbf{T568B}：白橙、橙、白绿、蓝、白蓝、绿、白棕、棕
\end{itemize}

\subsection{线缆类型}

\begin{keybox}
\textbf{直连线}（两端相同标准）：连接\textbf{不同类型}设备
\begin{itemize}[itemsep=0pt]
    \item 计算机 ↔ 交换机
    \item 计算机 ↔ 路由器
    \item 交换机 ↔ 路由器
\end{itemize}

\textbf{交叉线}（一端T568A，一端T568B）：连接\textbf{相同类型}设备
\begin{itemize}[itemsep=0pt]
    \item 计算机 ↔ 计算机
    \item 交换机 ↔ 交换机
    \item 路由器 ↔ 路由器
\end{itemize}
\end{keybox}

\subsection{双绞线使用的线对}

\begin{keybox}
10/100Mbps以太网只使用\textbf{1、2、3、6}号线（两对）。

因此，如果4、5、7、8号线不通，网线\textbf{仍可使用}。

如果1、2、3、6任一根不通，网线\textbf{不能使用}。

\textbf{双绞线最大长度}：100米（超过可能因\textbf{信号衰减}导致不能正常工作）
\end{keybox}

%========================================
\section{数据包分析}
%========================================

\subsection{协议层次}

\begin{keybox}
分析抓包数据时，识别各层协议：
\begin{itemize}[itemsep=0pt]
    \item \textbf{数据链路层}：Ethernet II
    \item \textbf{网络层}：IP、ICMP、ARP
    \item \textbf{传输层}：TCP、UDP
    \item \textbf{应用层}：HTTP、DNS、FTP等
\end{itemize}
\end{keybox}

\subsection{常见抓包分析}

\begin{exambox}[DNS查询分析]
执行\texttt{nslookup www.baidu.com}时抓包：
\begin{itemize}[itemsep=0pt]
    \item 应用层协议：DNS
    \item 传输层协议：UDP
    \item 目的端口：53
    \item 目的IP：DNS服务器地址
\end{itemize}
\end{exambox}

\begin{exambox}[Ping命令分析]
执行\texttt{ping 目标IP}时抓包：
\begin{itemize}[itemsep=0pt]
    \item 网络层协议：ICMP
    \item ICMP类型：Echo Request（请求）/ Echo Reply（回复）
\end{itemize}
\end{exambox}

%========================================
\section{综合题型总结}
%========================================

\subsection{VLAN配置题}

\begin{enumerate}[leftmargin=2em,itemsep=0pt]
    \item 创建VLAN并命名
    \item 将端口划入VLAN
    \item 配置Trunk端口（交换机互连）
    \item 配置VLAN间路由（三层交换机）
\end{enumerate}

\subsection{路由器配置题}

\begin{enumerate}[leftmargin=2em,itemsep=0pt]
    \item 配置接口IP地址
    \item 配置静态路由或默认路由
    \item 配置ACL（如需要）
    \item 配置NAT（如需要）
\end{enumerate}

\subsection{故障排查顺序}

\begin{enumerate}[leftmargin=2em,itemsep=0pt]
    \item 检查物理层（网线、接口）
    \item 检查IP配置（地址、掩码、网关）
    \item 检查路由配置
    \item 检查ACL配置
\end{enumerate}

%========================================
\section{链路聚合技术}
%========================================

\begin{keybox}
\textbf{链路聚合}（Link Aggregation）：将\textbf{多个物理链路}组成\textbf{一个逻辑链路}。

\textbf{作用}：
\begin{itemize}[itemsep=0pt]
    \item 增加带宽（多条链路同时传输）
    \item 提高可靠性（一条链路故障，其他链路继续工作）
    \item 负载均衡
\end{itemize}
\end{keybox}

%========================================
\section{路由器接口类型}
%========================================

\begin{keybox}
\textbf{路由器常见接口}：

\begin{itemize}[itemsep=0pt]
    \item \textbf{以太网口}（FastEthernet/GigabitEthernet）：连接局域网
    \item \textbf{串口}（Serial）：点对点连接，需要配置时钟速率（DCE端）
    \item \textbf{Console口}：本地配置管理
\end{itemize}

\textbf{串口连接}：
\begin{itemize}[itemsep=0pt]
    \item 串口之间必须\textbf{成对}连接（一端DCE，一端DTE）
    \item DCE端需要配置\texttt{clock rate}
    \item 串口IP地址通常在\textbf{同一网段}
\end{itemize}
\end{keybox}

%========================================
\section{网络故障排查}
%========================================

\subsection{常见故障现象与原因}

\begin{exambox}[故障1：无法访问域名]
\textbf{现象}：ping域名超时，但ping IP地址正常

\textbf{原因}：DNS解析故障

\textbf{解决}：
\begin{itemize}[itemsep=0pt]
    \item 检查DNS服务器配置
    \item 使用\texttt{nslookup}测试DNS解析
    \item 清空DNS缓存：\texttt{ipconfig /flushdns}
\end{itemize}
\end{exambox}

\begin{exambox}[故障2：无法访问外网]
\textbf{现象}：ping本网段正常，ping外网超时

\textbf{可能原因}：
\begin{itemize}[itemsep=0pt]
    \item 默认网关未设置或设置错误
    \item 路由器上未配置路由
    \item 路由器接口未启用（shutdown状态）
\end{itemize}
\end{exambox}

\begin{exambox}[故障3：Destination Host Unreachable]
\textbf{现象}：ping时显示"Destination Host Unreachable"

\textbf{原因}：
\begin{itemize}[itemsep=0pt]
    \item 本机到目标主机的路由不通
    \item 目标主机不存在或未开机
    \item 中间路由器无法转发
\end{itemize}
\end{exambox}

\subsection{故障排查工具}

\begin{itemize}[leftmargin=2em,itemsep=0pt]
    \item \textbf{ping}：测试连通性
    \item \textbf{tracert}：跟踪路由路径
    \item \textbf{ipconfig}：查看IP配置
    \item \textbf{nslookup}：测试DNS解析
    \item \textbf{arp -a}：查看ARP缓存
    \item \textbf{netstat}：查看网络连接状态
\end{itemize}

\subsection{测试连接的正确顺序}

\begin{cmdbox}
\textbf{测试网络连通性的步骤}：
\begin{lstlisting}
1. ping 127.0.0.1        # 测试TCP/IP协议栈
2. ping 本机IP           # 测试网卡
3. ping 默认网关         # 测试本网段
4. ping 外网IP           # 测试路由
5. ping 外网域名         # 测试DNS
\end{lstlisting}
\end{cmdbox}

%========================================
\section{以太网链路聚合}
%========================================

\begin{keybox}
\textbf{测试链路聚合}：

如果两台交换机之间有多条链路聚合：
\begin{itemize}[itemsep=0pt]
    \item 断开其中一条链路，网络\textbf{仍能正常通信}（其他链路继续工作）
    \item 所有链路同时工作，带宽叠加
\end{itemize}

\textbf{注意}：如果没有配置链路聚合，多条链路会形成环路，需要STP阻塞部分端口。
\end{keybox}

%========================================
\section{ARP缓存}
%========================================

\begin{keybox}
\textbf{ARP缓存类型}：

\begin{itemize}[itemsep=0pt]
    \item \textbf{dynamic}（动态）：通过ARP协议自动学习，有老化时间
    \item \textbf{static}（静态）：手动添加，不会老化
\end{itemize}

\textbf{查看ARP缓存}：
\begin{lstlisting}
arp -a
\end{lstlisting}

\textbf{手动添加静态ARP}：
\begin{lstlisting}
arp -s 192.168.5.7 00-52-18-45-a2-67
\end{lstlisting}
\end{keybox}

%========================================
\section{交换机端口状态}
%========================================

\begin{keybox}
\textbf{交换机端口默认状态}：

在缺省配置下，交换机的所有端口：
\begin{itemize}[itemsep=0pt]
    \item 处于\textbf{启用状态}（up）
    \item 属于\textbf{VLAN 1}（默认VLAN）
    \item 可以互相通信
\end{itemize}
\end{keybox}

%========================================
\section{网络设备功能对比}
%========================================

\begin{table}[htbp]
\centering\small
\begin{tabular}{lcccc}
\toprule
\textbf{功能} & \textbf{中继器} & \textbf{集线器} & \textbf{网桥} & \textbf{路由器} \\
\midrule
工作层次 & 物理层 & 物理层 & 数据链路层 & 网络层 \\
隔离冲突域 & 否 & 否 & 是 & 是 \\
隔离广播域 & 否 & 否 & 否 & 是 \\
转发依据 & — & — & MAC地址 & IP地址 \\
\bottomrule
\end{tabular}
\caption{网络设备功能对比}
\end{table}

\begin{keybox}
\textbf{抑制广播风暴}：只有\textbf{路由器}能隔离广播域，从而抑制广播风暴。
\end{keybox}

%========================================
\section{传输时延}
%========================================

\begin{keybox}
\textbf{传输时延}：信号从介质一端传到另一端所需的时间。

\textbf{影响因素}：
\begin{itemize}[itemsep=0pt]
    \item 介质长度
    \item 信号传播速度（光纤 > 双绞线）
\end{itemize}

\textbf{与传输时延相关的问题}：
\begin{itemize}[itemsep=0pt]
    \item \textbf{近端串扰}：信号在传输过程中相互干扰
    \item \textbf{衰减}：信号强度随距离增加而减弱
    \item \textbf{远端串扰}：远端的信号干扰
\end{itemize}
\end{keybox}

%========================================
\section{Windows网络组件}
%========================================

\begin{keybox}
\textbf{访问Web无关的组件}：\textbf{WINS}（Windows Internet Name Service）

WINS用于NetBIOS名称解析，与Web访问无关。

\textbf{Web访问相关组件}：
\begin{itemize}[itemsep=0pt]
    \item DNS：域名解析
    \item TCP：传输层协议
    \item IIS：Web服务器
\end{itemize}
\end{keybox}

%========================================
\section{路由器配置实验}
%========================================

\subsection{实验拓扑规划}

\begin{exambox}[典型实验拓扑]
\textbf{设备}：
\begin{itemize}[itemsep=0pt]
    \item 路由器RA：R3740（接口：s0/0, s0/1, s0/2, g0/0, g0/1）
    \item 路由器RB：R2692A（接口：s1/2, s1/3, f1/0, f1/1）
    \item 计算机PC1、PC2
\end{itemize}

\textbf{IP规划}（202.115.0.0-202.115.255.255，/24）：
\begin{itemize}[itemsep=0pt]
    \item PC1网段：202.115.10.0/24
    \item PC2网段：202.115.20.0/24
    \item 串口互连：202.115.30.0/24
\end{itemize}
\end{exambox}

\subsection{配置步骤}

\begin{cmdbox}
\textbf{1. 配置路由器接口}：
\begin{lstlisting}
RA(config)# interface gigabitEthernet 0/0
RA(config-if)# ip address 202.115.10.1 255.255.255.0
RA(config-if)# no shutdown
RA(config-if)# exit

RA(config)# interface serial 0/0
RA(config-if)# ip address 202.115.30.1 255.255.255.0
RA(config-if)# clock rate 64000
RA(config-if)# no shutdown
RA(config-if)# exit
\end{lstlisting}

\textbf{2. 配置静态路由}：
\begin{lstlisting}
RA(config)# ip route 202.115.20.0 255.255.255.0 202.115.30.2
\end{lstlisting}

\textbf{3. 配置计算机}：
\begin{itemize}[itemsep=0pt]
    \item PC1：IP=202.115.10.2，网关=202.115.10.1
    \item PC2：IP=202.115.20.2，网关=202.115.20.1
\end{itemize}
\end{cmdbox}

\subsection{测试串口连接}

\begin{keybox}
\textbf{测试串口是否接好}：

在路由器上执行：
\begin{lstlisting}
Router# show interfaces serial 0/0
\end{lstlisting}

查看接口状态：
\begin{itemize}[itemsep=0pt]
    \item \texttt{Serial0/0 is up, line protocol is up}：连接正常
    \item \texttt{Serial0/0 is up, line protocol is down}：物理连接正常，但协议未建立（可能时钟未配置）
    \item \texttt{Serial0/0 is down}：物理连接故障
\end{itemize}
\end{keybox}

%========================================
\section{VLAN实验}
%========================================

\subsection{实验要求}

\begin{exambox}[VLAN实验典型题]
\textbf{场景}：某单位有40台计算机，财务部4台，其他部门36台。使用两台交换机S3760和S2126。

\textbf{要求}：
\begin{itemize}[itemsep=0pt]
    \item 财务部划入VLAN2
    \item 其他部门划入VLAN3
    \item 使用私有IP地址10.0.0.0/8
    \item 财务部：10.1.1.0/24
    \item 其他部门：10.2.1.0/24
\end{itemize}
\end{exambox}

\subsection{配置命令}

\begin{cmdbox}
\textbf{S3760配置}：
\begin{lstlisting}
! 创建VLAN
S3760(config)# vlan 2
S3760(config-vlan)# name vlan2
S3760(config-vlan)# exit
S3760(config)# vlan 3
S3760(config-vlan)# name vlan3
S3760(config-vlan)# exit

! 指定端口属于VLAN
S3760(config)# interface fastEthernet 0/1
S3760(config-if)# switchport access vlan 2
S3760(config-if)# exit
S3760(config)# interface range fastEthernet 0/2-21
S3760(config-if-range)# switchport access vlan 3
S3760(config-if-range)# exit

! 配置Trunk端口
S3760(config)# interface fastEthernet 0/24
S3760(config-if)# switchport mode trunk
S3760(config-if)# exit
\end{lstlisting}

\textbf{配置VLAN间路由}（三层交换机）：
\begin{lstlisting}
! 进入VLAN 2的虚拟接口
S3760(config)# interface vlan 2
! 配置IP地址：10.1.1.254是网关地址，255.0.0.0是子网掩码
S3760(config-if)# ip address 10.1.1.254 255.0.0.0
S3760(config-if)# exit

! 进入VLAN 3的虚拟接口
S3760(config)# interface vlan 3
! 配置IP地址：10.2.1.254是网关地址，255.0.0.0是子网掩码
S3760(config-if)# ip address 10.2.1.254 255.0.0.0
S3760(config-if)# exit

! 启用三层交换机的路由功能，实现VLAN间通信
S3760(config)# ip routing
\end{lstlisting}
\end{cmdbox}

%========================================
\section{访问控制列表实验}
%========================================

\subsection{ACL配置实例}

\begin{exambox}[ACL实验题]
\textbf{要求}：使部门B不能访问部门A，其他访问不受限。

\textbf{假设}：
\begin{itemize}[itemsep=0pt]
    \item 部门A：172.16.12.0/24
    \item 部门B：172.16.13.0/24
\end{itemize}
\end{exambox}

\begin{cmdbox}
\textbf{配置扩展ACL}：
\begin{lstlisting}
! 创建ACL规则
Router(config)# access-list 101 deny ip 172.16.13.0 0.0.0.255 172.16.12.0 0.0.0.255
Router(config)# access-list 101 permit ip any any

! 应用到接口（部门B的接口）
Router(config)# interface fastEthernet 0/1
Router(config-if)# ip access-group 101 in
Router(config-if)# exit
\end{lstlisting}
\end{cmdbox}

\subsection{限制FTP访问}

\begin{exambox}[限制FTP访问]
\textbf{要求}：禁止内网访问外网的FTP服务器。
\end{exambox}

\begin{cmdbox}
\begin{lstlisting}
! FTP使用TCP端口21（控制）和20（数据）
Router(config)# access-list 102 deny tcp 192.168.1.0 0.0.0.255 any eq 21
Router(config)# access-list 102 deny tcp 192.168.1.0 0.0.0.255 any eq 20
Router(config)# access-list 102 permit ip any any

! 应用到内网接口的出方向
Router(config)# interface fastEthernet 0/0
Router(config-if)# ip access-group 102 out
\end{lstlisting}
\end{cmdbox}

%========================================
\section{NAT实验}
%========================================

\subsection{NAT配置完整示例}

\begin{exambox}[NAT实验]
\textbf{拓扑}：
\begin{itemize}[itemsep=0pt]
    \item 内网：192.168.1.0/24
    \item 路由器内部接口：f0/0 = 192.168.1.1
    \item 路由器外部接口：s0/0 = 59.67.148.3
    \item 外网服务器：202.113.64.2
\end{itemize}
\end{exambox}

\begin{cmdbox}
\textbf{完整NAT配置}：
\begin{lstlisting}
! 定义允许转换的内网地址
Router(config)# access-list 11 permit 192.168.1.0 0.0.0.255

! 定义内部全局地址池
Router(config)# ip nat pool abc 59.67.148.3 59.67.148.3 netmask 255.255.255.0

! 配置动态NAT（PAT）
Router(config)# ip nat inside source list 11 pool abc overload

! 指定内部接口
Router(config)# interface fastEthernet 0/0
Router(config-if)# ip nat inside
Router(config-if)# exit

! 指定外部接口
Router(config)# interface serial 0/0
Router(config-if)# ip nat outside
Router(config-if)# exit
\end{lstlisting}

\textbf{查看NAT转换表}：
\begin{lstlisting}
Router# show ip nat translations
\end{lstlisting}
\end{cmdbox}

%========================================
\section{数据包分析实验}
%========================================

\subsection{Wireshark抓包分析}

\begin{exambox}[DNS查询抓包]
\textbf{执行命令}：\texttt{nslookup www.baidu.com}

\textbf{抓包分析}：
\begin{itemize}[itemsep=0pt]
    \item 数据链路层：Ethernet II
    \item 网络层：IP
    \item 传输层：UDP
    \item 应用层：DNS
    \item 源IP：本机IP（如118.114.175.86）
    \item 目的IP：DNS服务器IP（如61.139.2.69）
    \item 源端口：随机高端口
    \item 目的端口：53
    \item 源MAC：本机MAC
    \item 目的MAC：网关MAC
\end{itemize}
\end{exambox}

\begin{exambox}[HTTP访问抓包]
\textbf{执行命令}：访问网页

\textbf{抓包分析}：
\begin{itemize}[itemsep=0pt]
    \item 应用层：HTTP
    \item 传输层：TCP
    \item 目的端口：80（HTTP）或443（HTTPS）
    \item 包含三次握手过程
\end{itemize}
\end{exambox}

\subsection{TCP连接分析}

\begin{keybox}
\textbf{TCP三次握手抓包特征}：

\textbf{第一次握手}：
\begin{itemize}[itemsep=0pt]
    \item SYN=1, ACK=0
    \item Sequence Number=x（随机初始值）
\end{itemize}

\textbf{第二次握手}：
\begin{itemize}[itemsep=0pt]
    \item SYN=1, ACK=1
    \item Sequence Number=y（服务器随机初始值）
    \item Acknowledgement Number=x+1
\end{itemize}

\textbf{第三次握手}：
\begin{itemize}[itemsep=0pt]
    \item SYN=0, ACK=1
    \item Sequence Number=x+1
    \item Acknowledgement Number=y+1
\end{itemize}
\end{keybox}

%========================================
\section{综合实验题}
%========================================

\subsection{完整网络配置实验}

\begin{exambox}[综合实验]
\textbf{场景}：某单位内部局域网通过路由器连接Internet。

\textbf{要求}：
\begin{enumerate}[itemsep=0pt]
    \item 配置VLAN划分部门
    \item 配置路由器实现内网互通
    \item 配置NAT实现内网访问外网
    \item 配置ACL限制FTP访问
\end{enumerate}

\textbf{配置步骤}：
\begin{enumerate}[itemsep=0pt]
    \item 交换机：创建VLAN，划分端口，配置Trunk
    \item 三层交换机：配置VLAN接口IP，启用路由
    \item 路由器：配置接口IP，配置静态路由
    \item 路由器：配置NAT（ACL+地址池+inside/outside）
    \item 路由器：配置ACL限制访问
    \item 计算机：配置IP、掩码、网关、DNS
    \item 测试：ping、tracert、访问网页
\end{enumerate}
\end{exambox}


%========================================
\section{DHCP动态主机配置协议}
%========================================

\subsection{DHCP工作原理}

\begin{keybox}
\textbf{DHCP四步过程}（重要考点）：

\begin{enumerate}[itemsep=0pt]
    \item \textbf{DHCP Discover}：客户端广播寻找DHCP服务器
    \item \textbf{DHCP Offer}：服务器提供IP地址
    \item \textbf{DHCP Request}：客户端请求使用该IP
    \item \textbf{DHCP ACK}：服务器确认分配
\end{enumerate}

\textbf{DHCP使用UDP协议}：
\begin{itemize}[itemsep=0pt]
    \item 客户端端口：68
    \item 服务器端口：67
\end{itemize}
\end{keybox}

\subsection{DHCP配置}

\begin{cmdbox}
\textbf{路由器作为DHCP服务器}：
\begin{lstlisting}
! 创建DHCP地址池
Router(config)# ip dhcp pool mypool
Router(dhcp-config)# network 192.168.1.0 255.255.255.0
Router(dhcp-config)# default-router 192.168.1.1
Router(dhcp-config)# dns-server 8.8.8.8
Router(dhcp-config)# exit

! 排除不分配的地址
Router(config)# ip dhcp excluded-address 192.168.1.1 192.168.1.10
\end{lstlisting}
\end{cmdbox}

\begin{exambox}[DHCP抓包分析]
\textbf{执行}：计算机设置为自动获取IP

\textbf{抓包特征}：
\begin{itemize}[itemsep=0pt]
    \item 传输层协议：UDP
    \item 源端口：68，目的端口：67
    \item 源IP：0.0.0.0（初始状态）
    \item 目的IP：255.255.255.255（广播）
    \item 包含四个报文：Discover、Offer、Request、ACK
\end{itemize}
\end{exambox}

%========================================
\section{交换机MAC地址表}
%========================================

\subsection{MAC地址学习}

\begin{keybox}
\textbf{交换机自学习过程}：

\begin{enumerate}[itemsep=0pt]
    \item 交换机收到数据帧，记录\textbf{源MAC地址}和\textbf{进入端口}
    \item 查找MAC地址表中的\textbf{目的MAC地址}
    \item 如果找到，从对应端口转发
    \item 如果未找到，\textbf{泛洪}（除接收端口外的所有端口）
\end{enumerate}

\textbf{MAC地址表项类型}：
\begin{itemize}[itemsep=0pt]
    \item \textbf{动态}：自动学习，有老化时间（默认300秒）
    \item \textbf{静态}：手动配置，不会老化
\end{itemize}
\end{keybox}

\subsection{查看MAC地址表}

\begin{cmdbox}
\begin{lstlisting}
Switch# show mac-address-table
Switch# show mac-address-table dynamic
Switch# show mac-address-table static
\end{lstlisting}
\end{cmdbox}

%========================================
\section{端口镜像(Port Mirroring)}
%========================================

\begin{keybox}
\textbf{端口镜像作用}：将一个或多个端口的流量复制到监控端口，用于\textbf{网络监控和故障排查}。

\textbf{应用场景}：
\begin{itemize}[itemsep=0pt]
    \item 使用Wireshark抓包分析
    \item 网络安全监控
    \item 流量分析
\end{itemize}
\end{keybox}

\begin{cmdbox}
\textbf{配置端口镜像}：
\begin{lstlisting}
! 将f0/1的流量镜像到f0/24
Switch(config)# monitor session 1 source interface f0/1
Switch(config)# monitor session 1 destination interface f0/24
\end{lstlisting}
\end{cmdbox}

%========================================
\section{路由协议基础}
%========================================

\subsection{静态路由vs动态路由}

\begin{table}[htbp]
\centering
\begin{tabular}{lcc}
\toprule
\textbf{特性} & \textbf{静态路由} & \textbf{动态路由} \\
\midrule
配置方式 & 手动配置 & 自动学习 \\
适用场景 & 小型网络 & 大型网络 \\
网络变化 & 需手动更新 & 自动适应 \\
资源消耗 & 低 & 高 \\
管理难度 & 简单网络易，复杂网络难 & 配置后自动维护 \\
\bottomrule
\end{tabular}
\caption{静态路由与动态路由对比}
\end{table}

\subsection{RIP协议}

\begin{keybox}
\textbf{RIP特点}：
\begin{itemize}[itemsep=0pt]
    \item 距离矢量路由协议
    \item 以\textbf{跳数}作为度量（最大15跳，16跳表示不可达）
    \item 每30秒更新一次路由表
    \item 使用UDP端口520
\end{itemize}
\end{keybox}

\begin{cmdbox}
\textbf{配置RIP}：
\begin{lstlisting}
Router(config)# router rip
Router(config-router)# version 2
Router(config-router)# network 192.168.1.0
Router(config-router)# network 192.168.2.0
Router(config-router)# exit
\end{lstlisting}
\end{cmdbox}

\subsection{OSPF协议}

\begin{keybox}
\textbf{OSPF特点}：
\begin{itemize}[itemsep=0pt]
    \item 链路状态路由协议
    \item 以\textbf{带宽}作为度量（开销cost）
    \item 收敛速度快
    \item 支持大型网络
    \item 使用IP协议号89
\end{itemize}
\end{keybox}

\begin{cmdbox}
\textbf{配置OSPF}：
\begin{lstlisting}
Router(config)# router ospf 1
Router(config-router)# network 192.168.1.0 0.0.0.255 area 0
Router(config-router)# network 192.168.2.0 0.0.0.255 area 0
Router(config-router)# exit
\end{lstlisting}
\end{cmdbox}

%========================================
\section{网络安全基础}
%========================================

\subsection{端口安全}

\begin{keybox}
\textbf{端口安全作用}：限制交换机端口上允许的MAC地址，防止未授权设备接入。

\textbf{违规处理方式}：
\begin{itemize}[itemsep=0pt]
    \item \textbf{protect}：丢弃非法流量，不记录
    \item \textbf{restrict}：丢弃非法流量，记录日志
    \item \textbf{shutdown}：关闭端口（默认）
\end{itemize}
\end{keybox}

\begin{cmdbox}
\textbf{配置端口安全}：
\begin{lstlisting}
Switch(config)# interface f0/1
Switch(config-if)# switchport mode access
Switch(config-if)# switchport port-security
Switch(config-if)# switchport port-security maximum 2
Switch(config-if)# switchport port-security violation shutdown
Switch(config-if)# switchport port-security mac-address sticky
Switch(config-if)# exit
\end{lstlisting}
\end{cmdbox}

\subsection{常见网络攻击}

\begin{keybox}
\textbf{ARP欺骗}：攻击者发送伪造的ARP响应，将网关MAC地址映射到攻击者MAC。

\textbf{防御措施}：
\begin{itemize}[itemsep=0pt]
    \item 配置静态ARP绑定
    \item 使用ARP防火墙
    \item 启用端口安全
\end{itemize}

\textbf{MAC地址泛洪攻击}：攻击者发送大量不同源MAC的数据帧，填满交换机MAC地址表。

\textbf{防御措施}：
\begin{itemize}[itemsep=0pt]
    \item 配置端口安全
    \item 限制每端口MAC地址数量
\end{itemize}
\end{keybox}

%========================================
\section{网络性能测试}
%========================================

\subsection{带宽测试}

\begin{keybox}
\textbf{理论传输时间计算}：

传输时间 = 数据大小 / 带宽

\textbf{示例}：100MB文件通过100Mbps链路传输
\begin{itemize}[itemsep=0pt]
    \item 数据大小 = 100MB = 100 × 8 = 800Mb
    \item 传输时间 = 800Mb / 100Mbps = 8秒
\end{itemize}

\textbf{注意}：实际传输时间会因协议开销、网络拥塞等因素增加。
\end{keybox}

\subsection{时延类型}

\begin{keybox}
\textbf{网络时延组成}：

\begin{enumerate}[itemsep=0pt]
    \item \textbf{发送时延}：数据从主机发送到链路所需时间 = 数据长度 / 发送速率
    \item \textbf{传播时延}：信号在链路上传播所需时间 = 链路长度 / 传播速度
    \item \textbf{处理时延}：路由器处理分组所需时间
    \item \textbf{排队时延}：分组在路由器队列中等待的时间
\end{enumerate}

\textbf{总时延} = 发送时延 + 传播时延 + 处理时延 + 排队时延
\end{keybox}

%========================================
\section{IPv6基础}
%========================================

\subsection{IPv6地址格式}

\begin{keybox}
\textbf{IPv6地址}：128位，用冒号分隔的8组十六进制数。

\textbf{示例}：2001:0db8:85a3:0000:0000:8a2e:0370:7334

\textbf{简化规则}：
\begin{itemize}[itemsep=0pt]
    \item 每组前导0可省略：2001:db8:85a3:0:0:8a2e:370:7334
    \item 连续的0组可用::表示（只能用一次）：2001:db8:85a3::8a2e:370:7334
\end{itemize}
\end{keybox}

%========================================
\section{无线网络基础}
%========================================

\subsection{无线网络标准}

\begin{table}[htbp]
\centering
\begin{tabular}{lccc}
\toprule
\textbf{标准} & \textbf{频段} & \textbf{最大速率} & \textbf{特点} \\
\midrule
802.11b & 2.4GHz & 11Mbps & 早期标准 \\
802.11g & 2.4GHz & 54Mbps & 兼容11b \\
802.11n & 2.4/5GHz & 600Mbps & 双频，MIMO \\
802.11ac & 5GHz & 1.3Gbps & 高速 \\
\bottomrule
\end{tabular}
\caption{无线网络标准对比}
\end{table}

\subsection{无线安全}

\begin{keybox}
\textbf{无线加密方式}（安全性从低到高）：

\begin{enumerate}[itemsep=0pt]
    \item \textbf{开放}：无加密（不安全）
    \item \textbf{WEP}：已被破解（不推荐）
    \item \textbf{WPA}：较安全
    \item \textbf{WPA2}：推荐使用
    \item \textbf{WPA3}：最新标准
\end{enumerate}

\textbf{认证方式}：
\begin{itemize}[itemsep=0pt]
    \item \textbf{PSK}（预共享密钥）：家庭/小型网络
    \item \textbf{802.1X}：企业级认证
\end{itemize}
\end{keybox}

%========================================
\section{网络管理协议}
%========================================

\subsection{SNMP简单网络管理协议}

\begin{keybox}
\textbf{SNMP组成}：
\begin{itemize}[itemsep=0pt]
    \item \textbf{管理站}：网络管理系统
    \item \textbf{代理}：运行在被管理设备上
    \item \textbf{MIB}：管理信息库
\end{itemize}

\textbf{SNMP版本}：
\begin{itemize}[itemsep=0pt]
    \item SNMPv1：基础版本
    \item SNMPv2c：增加批量操作
    \item SNMPv3：增加安全认证
\end{itemize}

\textbf{使用UDP端口161（代理）和162（trap）}
\end{keybox}

\subsection{Syslog日志}

\begin{keybox}
\textbf{Syslog严重级别}（0-7）：

\begin{itemize}[itemsep=0pt]
    \item 0 - Emergency：系统不可用
    \item 1 - Alert：需要立即处理
    \item 2 - Critical：严重错误
    \item 3 - Error：错误
    \item 4 - Warning：警告
    \item 5 - Notice：通知
    \item 6 - Informational：信息
    \item 7 - Debug：调试
\end{itemize}

\textbf{使用UDP端口514}
\end{keybox}

\subsection{易错点总结}

\begin{keybox}
\textbf{常见错误}：

\begin{enumerate}[itemsep=0pt]
    \item 混淆集线器和交换机的工作层次
    \item 认为VLAN可以自动配置
    \item 忘记配置默认网关导致无法跨网段通信
    \item ACL规则顺序错误（应该先具体后宽泛）
    \item NAT的inside/outside接口配置反了
    \item 串口忘记配置clock rate（DCE端）
    \item 接口配置后忘记no shutdown
    \item 子网掩码计算错误（注意/24、/16等表示法）
    \item 混淆源端口和目的端口
    \item 认为交换机能隔离广播域（只有路由器能）
\end{enumerate}
\end{keybox}

%========================================
\section{交换机配置模式详解}
%========================================

\subsection{交换机配置模式}

交换机命令行操作模式主要包括：用户模式、特权模式、全局配置模式和端口配置模式等几种。

\begin{table}[htbp]
\centering
\begin{tabular}{ll}
\toprule
\textbf{模式} & \textbf{提示符} \\
\midrule
用户模式 & Switch> \\
特权模式 & Switch\# \\
全局配置模式 & Switch(config)\# \\
端口配置模式 & Switch(config-if)\# \\
VLAN配置模式 & Switch(config-vlan)\# \\
\bottomrule
\end{tabular}
\caption{交换机配置模式}
\end{table}

\subsection{模式切换命令}

\begin{cmdbox}
\textbf{用户模式进入特权模式}：
\begin{lstlisting}
Switch> enable
Password:
\end{lstlisting}
密码输入时在屏幕上不显示。

\textbf{特权模式进入全局配置模式}：
\begin{lstlisting}
Switch# configure terminal
Switch(config)# exit
\end{lstlisting}

\textbf{全局模式进入端口配置模式}：
\begin{lstlisting}
Switch(config)# interface port-type mod-num/port-num
! port-type: 接口类型(fastEthernet/gigabitEthernet/serial等)
! mod-num: 接口模块号
! port-num: 接口号
Switch(config-if)# exit
\end{lstlisting}

\textbf{示例}：
\begin{lstlisting}
Switch(config)# interface fastEthernet 0/1
! 进入快速以太网接口0/1的配置模式
\end{lstlisting}
\end{cmdbox}

\subsection{获得帮助}

\begin{cmdbox}
\begin{lstlisting}
Switch# ?                    ! 显示当前模式下所有命令
Switch# show ?               ! 显示所有以show开头的命令
\end{lstlisting}
\end{cmdbox}


%========================================
\section{生成树协议(STP)配置实验}
%========================================

\subsection{实验目的}

\begin{itemize}[leftmargin=2em,itemsep=0pt]
    \item 理解生成树协议的作用
    \item 掌握生成树协议的配置方法
    \item 理解根交换机和根端口的选举规则
\end{itemize}

\subsection{实验背景}

当交换机之间有多条链路连接时，会形成环路，造成广播风暴。生成树协议可以消除环路，使只有一条线路处于工作状态。

\subsection{配置步骤}

\begin{cmdbox}
\textbf{1. 在交换机Sa上开启生成树协议}：
\begin{lstlisting}
Sa(config)# spanning-tree
! spanning-tree: 启用生成树协议
! 交换机默认情况下生成树协议没有工作
\end{lstlisting}

\textbf{停止生成树协议}：
\begin{lstlisting}
Sa(config)# no spanning-tree
! no spanning-tree: 关闭生成树协议
\end{lstlisting}

\textbf{2. 设置生成树协议类型}：
\begin{lstlisting}
Sa(config)# spanning-tree mode rstp
! spanning-tree mode: 设置生成树协议类型
! rstp: 快速生成树协议(Rapid Spanning Tree Protocol)
! 其他选项: stp(标准生成树), mstp(多生成树)
\end{lstlisting}

\textbf{3. 设置交换机Sa的优先级}（使它成为根交换机）：
\begin{lstlisting}
Sa(config)# spanning-tree priority 4096
! spanning-tree priority: 设置交换机优先级
! 4096: 优先级值(必须是4096的倍数)
! 范围: 0-61440, 默认值32768
! 数值越小,优先级越高,越可能成为根交换机
\end{lstlisting}

\textbf{4. 设置端口优先级}：
\begin{lstlisting}
Sa(config)# interface fastEthernet 0/3
Sa(config-if)# spanning-tree port-priority 0
! spanning-tree port-priority: 设置端口优先级
! 0: 端口优先级值(必须是16的倍数)
! 范围: 0-240
! 数值越小,优先级越高,越可能成为根端口
\end{lstlisting}

\textbf{5. 查看生成树配置信息}：
\begin{lstlisting}
Sa# show spanning-tree
! show spanning-tree: 显示生成树协议状态和配置
! 包括: 根桥ID、端口状态、优先级等信息
\end{lstlisting}
\end{cmdbox}

\begin{keybox}
\textbf{生成树选举规则}：

\textbf{根交换机选举}：
\begin{itemize}[itemsep=0pt]
    \item 选择最小的交换机优先级
    \item 优先级相同则选最小MAC地址
\end{itemize}

\textbf{根端口选举}：
\begin{itemize}[itemsep=0pt]
    \item 非根交换机上选择到根交换机路径开销最小的端口
    \item 开销相同则选最小端口编号
\end{itemize}

\textbf{重要参数}：
\begin{itemize}[itemsep=0pt]
    \item Priority：优先级
    \item RootCost：到根交换机的开销
    \item RootPort：根端口
    \item DesignatedRoot：根交换机ID
\end{itemize}
\end{keybox}


%========================================
\section{链路聚合配置实验}
%========================================

\subsection{实验目的}

\begin{itemize}[leftmargin=2em,itemsep=0pt]
    \item 理解链路聚合的作用
    \item 掌握链路聚合的配置方法
    \item 理解流量均衡的配置
\end{itemize}

\subsection{配置前准备}

配置链路聚合前要先删除生成树配置。默认情况下，交换机中不存在AggregatePort(AP)。

\subsection{配置步骤}

\begin{cmdbox}
\textbf{1. 在交换机Sa创建链路聚合接口AP1}：
\begin{lstlisting}
Sa(config)# interface aggregatePort 1
! interface aggregatePort: 创建聚合端口(AP)
! 1: AP编号(锐捷交换机支持1-6)
Sa(config-if)# switchport mode trunk
! switchport mode trunk: 设置AP工作模式为trunk
! trunk模式允许多个VLAN的数据通过
Sa(config-if)# exit
\end{lstlisting}

\textbf{2. 将网络接口F0/3和F0/4加入AP1}：
\begin{lstlisting}
Sa(config)# interface fastEthernet 0/3
Sa(config-if)# port-group 1
! port-group: 将接口加入指定的聚合端口组
! 1: AP编号,与上面创建的AP1对应
Sa(config)# interface fastEthernet 0/4
Sa(config-if)# port-group 1
! 将F0/4也加入AP1
! 注意: 加入AP后,接口属性将被AP属性覆盖
\end{lstlisting}

\textbf{3. 从AP中删除接口}：
\begin{lstlisting}
Sa(config-if)# no port-group
! no port-group: 将接口从聚合端口组中移除
! 移除后接口恢复加入AP前的属性
\end{lstlisting}

\textbf{4. 在交换机Sb上进行相同配置}：
\begin{lstlisting}
Sb(config)# interface aggregatePort 1
Sb(config-if)# switchport mode trunk
Sb(config-if)# exit
Sb(config)# interface fastEthernet 0/3
Sb(config-if)# port-group 1
Sb(config)# interface fastEthernet 0/4
Sb(config-if)# port-group 1
\end{lstlisting}

\textbf{5. 查看链路聚合配置信息}：
\begin{lstlisting}
Sb# show aggregatePort summary
! show aggregatePort summary: 显示聚合端口摘要信息
! 包括: 端口模式、成员端口、最大端口数等
\end{lstlisting}

\textbf{6. 配置聚合端口的流量均衡}：
\begin{lstlisting}
Sa# aggregateport load-balance dst-mac
! aggregateport load-balance: 配置流量均衡算法
! dst-mac: 按目的MAC地址分配流量
! 其他选项: src-mac(源MAC), src-ip(源IP)
! 默认为src-mac(源MAC地址)
Sa# aggregateport load-balance src-ip
! src-ip: 按源IP地址分配流量
\end{lstlisting}

\textbf{7. 查看流量均衡配置}：
\begin{lstlisting}
Sa# show aggregateport load-balance
! 显示当前使用的流量均衡算法
\end{lstlisting}
\end{cmdbox}

\begin{keybox}
\textbf{链路聚合测试}：

将交换机间的两条连线都连上，在PC1上执行：
\begin{lstlisting}
C:\> ping -t 192.168.10.12
\end{lstlisting}

拔掉F0/3口的连线，PC1与PC3之间仍然连通，一般没有丢包的情况。这说明链路聚合配置成功。
\end{keybox}


%========================================
\section{路由器基本配置实验}
%========================================

\subsection{路由器命令模式}

\begin{table}[htbp]
\centering
\begin{tabular}{lll}
\toprule
\textbf{命令模式} & \textbf{路由器提示符} & \textbf{访问/退出方法} \\
\midrule
用户EXEC & R2600-A> & 登录/使用exit命令 \\
特权EXEC & R2600-A\# & enable/disable或exit \\
全局配置 & R2600-A(config)\# & configure terminal/exit或end \\
接口配置 & R2600-A(config-if)\# & interface 接口/exit或end \\
\bottomrule
\end{tabular}
\caption{路由器命令模式及切换}
\end{table}

\subsection{路由器命令行基本功能}

\begin{cmdbox}
\textbf{1. 帮助信息}：
\begin{lstlisting}
R2600-A> ?                   ! 显示当前模式下所有可执行的命令
R2600-A# co?                 ! 显示所有以co开头的命令
R2600-A# configure ?         ! 显示configure命令后可执行的参数
\end{lstlisting}

\textbf{2. 命令的简写}：
\begin{lstlisting}
R2600-A# conf t              ! 代表configure terminal
\end{lstlisting}
只要和其他命令不引起歧义即可。

\textbf{3. 命令的自动补齐}：
\begin{lstlisting}
R2600-A# com<Tab>            ! 按Tab键自动补齐为configure
\end{lstlisting}

\textbf{4. 命令的快捷键功能}：
\begin{lstlisting}
Ctrl+Z                       ! 直接退回到特权模式
↑键                          ! 重新输入上一条命令
\end{lstlisting}
\end{cmdbox}

\subsection{路由器基本配置}

\begin{cmdbox}
\textbf{1. 配置设备名称}：
\begin{lstlisting}
R2600-A> enable
R2600-A# configure terminal
R2600-A(config)# hostname RA
\end{lstlisting}
注意：设备名称的有效字符是22个字节。

\textbf{2. 配置以太网接口IP地址}：
\begin{lstlisting}
RA(config)# interface fastEthernet 1/0
RA(config-if)# ip address 17.1.0.1 255.255.0.0
RA(config-if)# no shutdown
RA(config-if)# exit
\end{lstlisting}

\textbf{3. 配置串口IP地址}：
\begin{lstlisting}
RA(config)# interface serial 1/3
RA(config-if)# ip address 17.2.0.1 255.255.0.0
RA(config-if)# clock rate 64000    ! DCE端需要配置时钟速率
RA(config-if)# no shutdown
RA(config-if)# exit
\end{lstlisting}

\textbf{4. 查看路由器接口参数}：
\begin{lstlisting}
RA# show ip interface brief
\end{lstlisting}

\textbf{5. 查看当前生效的配置}：
\begin{lstlisting}
RA# show running-config
\end{lstlisting}

\textbf{6. 查看保存在NVRAM里的配置}：
\begin{lstlisting}
RA# show startup-config
\end{lstlisting}

\textbf{7. 保存配置}：
\begin{lstlisting}
RA# copy running-config startup-config
\end{lstlisting}
\end{cmdbox}

\begin{keybox}
\textbf{串口连接注意事项}：

\begin{itemize}[itemsep=0pt]
    \item 串口之间必须成对连接（一端DCE，一端DTE）
    \item DCE端需要配置clock rate
    \item 串口IP地址通常在同一网段
    \item 使用show interface命令查看接口状态
\end{itemize}

\textbf{接口状态说明}：
\begin{itemize}[itemsep=0pt]
    \item Serial is up, line protocol is up：连接正常
    \item Serial is up, line protocol is down：物理连接正常，但协议未建立
    \item Serial is down：物理连接故障
\end{itemize}
\end{keybox}


%========================================
\section{PPP协议配置}
%========================================

\subsection{PPP协议简介}

PPP协议分为LCP和NCP两个子层：
\begin{itemize}[leftmargin=2em,itemsep=0pt]
    \item \textbf{LCP}：负责链路的协商、建立、回拨、认证、数据的压缩、多链路捆绑等功能
    \item \textbf{NCP}：主要负责和上层协议进行协商，为网络层协议提供服务
\end{itemize}

\subsection{PPP配置命令}

\begin{cmdbox}
\textbf{配置广域网接口的封装类型为PPP}：
\begin{lstlisting}
Router(config)# interface serial 1/3
Router(config-if)# encapsulation ppp
Router(config-if)# exit
\end{lstlisting}
\end{cmdbox}

%========================================
\section{NAPT配置实验}
%========================================

\subsection{实验拓扑}

\begin{itemize}[leftmargin=2em,itemsep=0pt]
    \item 内网：10.0.0.0/16
    \item LAN-R的F1/0：10.0.0.1（内网接口）
    \item LAN-R的S1/3：202.115.12.1（外网接口）
    \item Internet-R的S1/3：202.115.12.2
    \item Internet-R的F1/0：200.2.0.1
    \item Web Server：200.2.0.2
    \item PC：10.0.0.2
\end{itemize}

\subsection{实验目的}

\begin{itemize}[leftmargin=2em,itemsep=0pt]
    \item 理解NAT/NAPT的工作原理
    \item 掌握NAT配置方法
    \item 理解内部本地地址和内部全局地址的概念
\end{itemize}

\subsection{配置步骤}

\begin{cmdbox}
\textbf{1. 在LAN-R上配置缺省路由}：
\begin{lstlisting}
LAN-R(config)# ip route 0.0.0.0 0.0.0.0 serial 1/3
\end{lstlisting}

\textbf{2. 在LAN-R上配置NAPT}：
\begin{lstlisting}
! 定义F1/0为内网接口
LAN-R(config)# interface fastEthernet 1/0
LAN-R(config-if)# ip nat inside
LAN-R(config-if)# exit

! 定义S1/3为外网接口
LAN-R(config)# interface serial 1/3
LAN-R(config-if)# ip nat outside
LAN-R(config-if)# exit

! 定义内部全局地址池
LAN-R(config)# ip nat pool to_internet 202.115.12.1 202.115.12.1 netmask 255.255.255.0

! 定义允许转换的地址
LAN-R(config)# access-list 10 permit 10.0.0.0 0.0.255.255

! 为内部本地地址转换地址池（overload表示PAT）
LAN-R(config)# ip nat inside source list 10 pool to_internet overload
\end{lstlisting}

\textbf{3. 验证测试}：
\begin{lstlisting}
! 在PC1上测试访问Web服务器
! 在路由器LAN-R上查看NAPT映射关系
LAN-R# show ip nat translation
\end{lstlisting}

输出结果示例：
\begin{lstlisting}
Outside global  Pro  Inside global      Inside local       Outside local
200.2.0.2:80    tcp  202.115.12.1:4924  10.0.0.2:4924     200.2.0.2:80
\end{lstlisting}

结果显示，10.0.0.2访问200.2.0.2的80端口，路由器将私有地址10.0.0.2转换为合法的外网地址202.115.12.1。
\end{cmdbox}


%========================================
\section{访问控制列表(ACL)扩展实验}
%========================================

\subsection{标准IP访问控制列表配置}

\begin{exambox}[实验要求]
配置标准IP访问控制列表，禁止202.115.12.0网络访问Web Server主机。
\end{exambox}

\begin{cmdbox}
\textbf{1. 在Internet-R上配置标准IP访问控制列表}：
\begin{lstlisting}
! 网络掩码是反掩码
Internet-R(config)# access-list 1 deny 202.115.12.0 0.0.0.255
! 其他网络互访不受限制
Internet-R(config)# access-list 1 permit any
\end{lstlisting}

\textbf{2. 把访问控制列表应用在接口上}：
\begin{lstlisting}
! 应用在尽量靠近目的地址的接口
Internet-R(config)# interface fastEthernet 1/0
! 检查该接口上方向为out的数据流量
Internet-R(config-if)# ip access-group 1 out
Internet-R(config-if)# end
\end{lstlisting}

\textbf{3. 查看访问控制列表}：
\begin{lstlisting}
Internet-R# show access-list 1
\end{lstlisting}

\textbf{4. 查看接口配置信息}：
\begin{lstlisting}
Internet-R# show ip interface fastEthernet 1/0
\end{lstlisting}

\textbf{5. 删除访问控制列表}：
\begin{lstlisting}
Internet-R(config)# interface fastEthernet 1/0
Internet-R(config-if)# no ip access-group 1 out
\end{lstlisting}
\end{cmdbox}

\subsection{扩展IP访问控制列表配置}

\begin{exambox}[实验要求]
配置扩展IP访问控制列表，禁止202.115.12.0网络访问Web Server主机的Web服务，其他网络互访不受限制。
\end{exambox}

\begin{cmdbox}
\textbf{1. 在Internet-R上配置扩展IP访问控制列表}：
\begin{lstlisting}
! eq www指示目的端口为80
Internet-R(config)# access-list 101 deny tcp 202.115.12.0 0.0.0.255 host 200.2.0.2 eq www
! 网络掩码是反掩码
Internet-R(config)# access-list 101 permit ip any any
\end{lstlisting}

\textbf{2. 把访问控制列表应用在接口上}：
\begin{lstlisting}
! 应用在尽量靠近目的地址的接口
Internet-R(config)# interface fastEthernet 1/0
! 检查该接口上方向为out的数据流量
Internet-R(config-if)# ip access-group 101 out
Internet-R(config-if)# end
\end{lstlisting}

\textbf{3. 验证测试}：

PC1不能访问Web Server的主页，但能PING通Web Server。
\end{cmdbox}

%========================================
\section{三层交换机VLAN间路由实验}
%========================================

\subsection{实验目的}

使分别属于两个不同静态VLAN（Vlan10和Vlan20）的计算机能够相互通信。

\subsection{配置步骤}

\begin{cmdbox}
\textbf{1. 配置VLAN}（参考前面VLAN配置实验）

\textbf{2. 配置Trunk}（参考前面Trunk配置）

\textbf{3. 配置三层交换机实现VLAN间通信}：

在三层交换机上创建虚拟VLAN接口，配置其IP，使不同VLAN间的计算机能够通信。

\begin{lstlisting}
! 创建虚拟Vlan10接口并配置其IP
S3760(config)# interface vlan 10
S3760(config-if)# ip address 192.168.10.1 255.255.255.0
S3760(config-if)# exit

! 创建虚拟Vlan20接口并配置其IP
S3760(config)# interface vlan 20
S3760(config-if)# ip address 192.168.20.1 255.255.255.0
S3760(config-if)# exit
\end{lstlisting}

\textbf{注意}：给VLAN配置的IP地址就是这个网段的网关地址，在该网段中计算机的网关地址应该设置为这个地址。

\textbf{4. 启动路由选择协议}：

在三层交换机上需要激活路由协议。只有这样当目的地址在本地VLAN上时，三层交换机才会使用路由协议来转发分组。

\begin{lstlisting}
S3760# configure terminal
S3760(config)# ip routing        ! 启动IP路由选择
S3760(config)# end
\end{lstlisting}

这一命令通常在默认情况下是启动的，生成的路由表也与路由器中基本相同。

\textbf{5. 查看配置结果}：
\begin{lstlisting}
S3760# show running-config
\end{lstlisting}

\textbf{6. 验证实验结果}：

配置完成后，任意两台计算机之间都应能通信。在任意一台计算机上执行ping另外三台主机的命令。
\end{cmdbox}


%========================================
\section{RIPv2路由协议配置实验}
%========================================

\subsection{实验目的}

\begin{itemize}[leftmargin=2em,itemsep=0pt]
    \item 理解距离矢量路由协议的工作原理
    \item 掌握RIPv2路由协议的配置方法
    \item 理解路由表的生成过程
\end{itemize}

\subsection{RIP协议特点}

\begin{itemize}[leftmargin=2em,itemsep=0pt]
    \item 距离矢量路由协议
    \item 使用跳数作为度量值
    \item 最大跳数为15
    \item 每30秒发送一次路由更新
    \item RIPv2支持VLSM和CIDR
\end{itemize}

\subsection{配置步骤}

\begin{cmdbox}
\textbf{1. 配置路由器接口}（参考前面路由器基本配置）

\textbf{2. 在路由器r2632-a上启用RIPv2协议}：
\begin{lstlisting}
r2632-a> enable
r2632-a# configure terminal
r2632-a(config)# router rip
! router rip: 进入RIP路由协议配置模式
r2632-a(config-router)# network 192.168.10.0
! network: 指定要启用RIP协议的网络
! 192.168.10.0: 网络地址(使用主类网络地址)
! RIP将在该网络的所有接口上发送和接收RIP更新
r2632-a(config-router)# network 192.168.12.0
r2632-a(config-router)# network 192.168.1.0
r2632-a(config-router)# version 2
! version 2: 指定使用RIPv2协议
! 默认为RIPv1, RIPv2支持VLSM和CIDR
\end{lstlisting}

\textbf{3. 在路由器r3740上启用RIPv2协议}：
\begin{lstlisting}
r3740> enable
r3740# configure terminal
r3740(config)# router rip
r3740(config-router)# network 192.168.10.0
r3740(config-router)# network 192.168.13.0
r3740(config-router)# network 192.168.0.0
r3740(config-router)# version 2
\end{lstlisting}

\textbf{4. 在其他路由器上进行类似配置}

\textbf{5. 查看路由表}：
\begin{lstlisting}
r3740# show ip route
\end{lstlisting}

输出结果说明：
\begin{itemize}[itemsep=0pt]
    \item C代表直连路由
    \item R表示通过RIP协议学到的可达的路由
    \item Via后面的IP地址表示要到达可达的网络经过的端口的IP地址
    \item [120/1]表示管理距离为120，跳数为1
\end{itemize}

\textbf{6. 查看RIP协议运行状态}：
\begin{lstlisting}
r3740# show ip protocols
\end{lstlisting}
\end{cmdbox}

\begin{keybox}
\textbf{RIPv2协议特点}：

\begin{itemize}[itemsep=0pt]
    \item 距离矢量路由协议
    \item 使用跳数作为度量值，最大跳数为15
    \item 每30秒发送一次路由更新
    \item 支持VLSM和CIDR
    \item 管理距离为120
\end{itemize}
\end{keybox}

%========================================
\section{OSPF路由协议配置实验}
%========================================

\subsection{实验目的}

\begin{itemize}[leftmargin=2em,itemsep=0pt]
    \item 理解链路状态路由协议的工作原理
    \item 掌握OSPF路由协议的配置方法
    \item 理解OSPF区域划分的概念
\end{itemize}

\subsection{OSPF协议特点}

\begin{itemize}[leftmargin=2em,itemsep=0pt]
    \item 链路状态路由协议
    \item 使用开销作为度量值
    \item 支持区域划分
    \item 收敛速度快
    \item 无路由环路问题
\end{itemize}

\subsection{OSPF区域划分}

OSPF协议支持区域划分，以减少资源消耗。本实验中将构建一个骨干区域Area 0，并把网络划分为多个区域。

\subsection{配置步骤}

\begin{cmdbox}
\textbf{1. 配置路由器接口}（参考前面路由器基本配置）

\textbf{2. 在r2632-b上进行Area 0区域OSPF配置}：
\begin{lstlisting}
r2632-b(config)# router ospf
! router ospf: 进入OSPF路由协议配置模式
r2632-b(config-router)# network 192.168.12.0 0.0.0.255 area 0
! network: 指定启用OSPF的网络
! 192.168.12.0: 网络地址
! 0.0.0.255: 反掩码(wildcard mask), 0表示必须匹配,1表示任意
! area 0: 指定区域号(0为骨干区域,所有区域必须连接到骨干区域)
r2632-b(config-router)# network 192.168.13.0 0.0.0.255 area 0
\end{lstlisting}

\textbf{3. 在r2632-a上进行Area 0和Area 1的配置}：
\begin{lstlisting}
r2632-a(config)# router ospf
! 进入OSPF协议配置模式
r2632-a(config-router)# network 192.168.12.0 0.0.0.255 area 0
! 将192.168.12.0网络划入骨干区域Area 0
r2632-a(config-router)# network 192.168.10.0 0.0.0.255 area 0
r2632-a(config-router)# network 192.168.1.0 0.0.0.255 area 1
! 将192.168.1.0网络划入Area 1
! 此路由器连接多个区域,成为ABR(区域边界路由器)
\end{lstlisting}

\textbf{4. 在r3740上进行Area 0和Area 2的配置}：
\begin{lstlisting}
r3740(config)# router ospf
r3740(config-router)# network 192.168.10.0 0.0.0.255 area 0
r3740(config-router)# network 192.168.13.0 0.0.0.255 area 0
r3740(config-router)# network 192.168.0.0 0.0.0.255 area 2
! 将192.168.0.0网络划入Area 2
\end{lstlisting}

\textbf{5. 在r1762上进行Area 2的配置}：
\begin{lstlisting}
r1762(config)# router ospf
r1762(config-router)# network 192.168.14.0 0.0.0.255 area 2
r1762(config-router)# network 192.168.0.0 0.0.0.255 area 2
! 所有接口都在Area 2中
\end{lstlisting}

\textbf{6. 查看OSPF协议运行状态}：
\begin{lstlisting}
r3740# show ip protocols
\end{lstlisting}

输出结果说明：
\begin{itemize}[itemsep=0pt]
    \item Router ID：路由器ID
    \item Number of areas：区域数量
    \item Routing for Networks：启用OSPF的网络
\end{itemize}

\textbf{7. 查看路由表}：
\begin{lstlisting}
r3740# show ip route
\end{lstlisting}

输出结果说明：
\begin{itemize}[itemsep=0pt]
    \item C代表直连路由
    \item O表示通过OSPF协议学到的路由
    \item O IA表示OSPF区域间路由
    \item [110/51]表示管理距离为110，开销为51
\end{itemize}

\textbf{8. 查看OSPF数据库}：
\begin{lstlisting}
r3740# show ip ospf database
\end{lstlisting}

\textbf{9. 验证主机连通性}：

在主机A上PING主机B，输出结果显示两台计算机可以连通。
\end{cmdbox}

\begin{keybox}
\textbf{OSPF协议特点}：

\begin{itemize}[itemsep=0pt]
    \item 链路状态路由协议
    \item 使用开销作为度量值
    \item 支持区域划分，减少资源消耗
    \item 收敛速度快
    \item 支持VLSM和CIDR
    \item 管理距离为110
    \item 使用Hello报文维护邻居关系
\end{itemize}

\textbf{OSPF报文类型}：
\begin{itemize}[itemsep=0pt]
    \item Hello报文：发现和维护邻居关系
    \item DBD报文：数据库描述报文
    \item LSR报文：链路状态请求报文
    \item LSU报文：链路状态更新报文
    \item LSAck报文：链路状态确认报文
\end{itemize}
\end{keybox}

%========================================
\section{路由表解读}
%========================================

\subsection{路由表标识符}

\begin{keybox}
\textbf{路由表中常见标识符}：
\begin{itemize}[itemsep=0pt]
    \item \textbf{C} - Connected，直连路由
    \item \textbf{S} - Static，静态路由
    \item \textbf{R} - RIP，通过RIP协议学习的路由
    \item \textbf{O} - OSPF，通过OSPF协议学习的路由
    \item \textbf{O IA} - OSPF Inter Area，OSPF区域间路由
    \item \textbf{B} - BGP，通过BGP协议学习的路由
    \item \textbf{*} - candidate default，候选默认路由
\end{itemize}
\end{keybox}

\subsection{路由表输出示例}

\begin{cmdbox}
\begin{lstlisting}
Router# show ip route
Codes: C-connected, S-static, R-RIP
       O-OSPF, IA-OSPF inter area
       
Gateway of last resort is not set

C    192.168.1.0/24 is directly connected, FastEthernet 1/0
! C表示直连路由，该网络直接连接在F1/0接口

S    192.168.2.0/24 [1/0] via 192.168.1.1
! S表示静态路由，[1/0]表示管理距离/度量值
! via 192.168.1.1表示下一跳地址

R    192.168.3.0/24 [120/1] via 192.168.1.2, 00:00:18, serial 1/3
! R表示RIP路由，[120/1]表示管理距离120，跳数1
! 00:00:18表示路由更新时间

O    192.168.4.0/24 [110/51] via 192.168.1.3, 00:03:50, serial 1/3
! O表示OSPF路由，[110/51]表示管理距离110，开销51
\end{lstlisting}
\end{cmdbox}

\subsection{管理距离对比}

\begin{table}[htbp]
\centering
\begin{tabular}{lc}
\toprule
\textbf{路由来源} & \textbf{管理距离} \\
\midrule
直连路由 & 0 \\
静态路由 & 1 \\
OSPF & 110 \\
RIP & 120 \\
\bottomrule
\end{tabular}
\caption{常见路由协议管理距离}
\end{table}

\begin{keybox}
\textbf{管理距离说明}：
\begin{itemize}[itemsep=0pt]
    \item 管理距离越小，路由优先级越高
    \item 当多个路由协议都有到达同一网络的路由时，选择管理距离最小的
    \item 直连路由管理距离为0，优先级最高
\end{itemize}
\end{keybox}

%========================================
\section{接口状态查看}
%========================================

\subsection{show ip interface brief}

\begin{cmdbox}
\begin{lstlisting}
Router# show ip interface brief
Interface         IP-Address      OK?  Status    Protocol
FastEthernet1/0   192.168.1.1     YES  UP        UP
FastEthernet1/1   unassigned      YES  DOWN      DOWN
Serial1/2         192.168.2.1     YES  UP        UP
Serial1/3         unassigned      YES  DOWN      DOWN
\end{lstlisting}

\textbf{字段说明}：
\begin{itemize}[itemsep=0pt]
    \item Interface: 接口名称
    \item IP-Address: 接口IP地址（unassigned表示未配置）
    \item OK?: 配置是否有效
    \item Status: 物理层状态（UP/DOWN）
    \item Protocol: 数据链路层协议状态（UP/DOWN）
\end{itemize}
\end{cmdbox}

\subsection{接口状态组合含义}

\begin{table}[htbp]
\centering
\begin{tabular}{lll}
\toprule
\textbf{Status} & \textbf{Protocol} & \textbf{含义} \\
\midrule
UP & UP & 接口正常工作 \\
UP & DOWN & 物理连接正常，协议未建立 \\
DOWN & DOWN & 物理连接故障或接口被shutdown \\
administratively down & DOWN & 接口被管理员手动关闭 \\
\bottomrule
\end{tabular}
\caption{接口状态组合}
\end{table}

%========================================
\section{show running-config详解}
%========================================

\begin{cmdbox}
\begin{lstlisting}
Router# show running-config
Building configuration...

Current configuration: 540 bytes
!
version 8.51
!
hostname RA                    ! 设备名称
!
no service password-encryption
!
interface Serial 1/2
 clock rate 64000              ! DCE端时钟速率
!
interface Serial 1/3
 ip address 17.2.0.1 255.255.0.0  ! 接口IP地址
!
interface FastEthernet 1/0
 ip address 17.1.0.1 255.255.0.0
 duplex auto                   ! 双工模式自动协商
 speed auto                    ! 速率自动协商
!
ip route 17.3.0.0 255.255.0.0 17.2.0.2  ! 静态路由
!
end
\end{lstlisting}
\end{cmdbox}

\begin{keybox}
\textbf{配置文件说明}：
\begin{itemize}[itemsep=0pt]
    \item \textbf{show running-config}：查看当前生效的配置（存储在RAM中）
    \item \textbf{show startup-config}：查看保存的配置（存储在NVRAM中）
    \item 路由器启动时，将NVRAM里的配置文件加载至RAM里生效
    \item 修改配置后需要保存：\texttt{copy running-config startup-config}
\end{itemize}
\end{keybox}

%========================================
\section{DCE与DTE配置}
%========================================

\begin{keybox}
\textbf{串口连接说明}：
\begin{itemize}[itemsep=0pt]
    \item \textbf{DCE}（Data Circuit-terminating Equipment）：数据电路终端设备，提供时钟信号
    \item \textbf{DTE}（Data Terminal Equipment）：数据终端设备，接收时钟信号
    \item 串口之间连接时，一端为DCE，一端为DTE
    \item \textbf{只有DCE端需要配置clock rate}
\end{itemize}

\textbf{判断DCE/DTE}：
\begin{lstlisting}
Router# show controllers serial 1/2
...
V35 DCE cable    ! 显示DCE表示该端为DCE
...
\end{lstlisting}

\textbf{配置DCE时钟速率}：
\begin{lstlisting}
Router(config)# interface serial 1/2
Router(config-if)# clock rate 64000
! clock rate: 时钟速率，单位bps
! 常用值: 64000, 128000, 256000, 512000, 1024000, 2048000
\end{lstlisting}
\end{keybox}

%========================================
\section{思考题汇总}
%========================================

\subsection{交换机相关}

\begin{enumerate}[leftmargin=2em,itemsep=0pt]
    \item 简述交换机的工作原理，与集线器相比，其优势在什么地方？
    \item 交换机是如何感知其周围网络环境变化的（如计算机增加、减少等）？
    \item 从一个网络管理员角度看，交换机提供了哪些功能供管理员方便、高效的管理所有网络用户？
    \item 什么是VLAN？该技术有哪些优点？如何实现VLAN？
    \item 什么是中继（Trunk）技术？如何实现Trunk？
\end{enumerate}

\subsection{路由器相关}

\begin{enumerate}[leftmargin=2em,itemsep=0pt]
    \item 简述路由器的工作原理，它对收到的数据包做哪些处理？
    \item 描述当用户访问Internet中某一服务器的过程中，源站发出的数据包是如何到达目的服务器的？中间经过了哪些封装与解封装，以及是如何选择路径的？
\end{enumerate}

\subsection{三层交换相关}

\begin{enumerate}[leftmargin=2em,itemsep=0pt]
    \item 说明三层交换技术产生的原因与技术特点
    \item 简述两种三层交换技术的工作原理
    \item 简述NHRP和MPLS各自的工作原理
\end{enumerate}

\subsection{路由协议相关}

\begin{enumerate}[leftmargin=2em,itemsep=0pt]
    \item 比较链路状态路由协议和距离矢量路由协议的优缺点
    \item 什么是管理距离AD（Administrative Distance）？OSPF的管理距离AD是多少？和RIP相比，有什么特点？
    \item 什么是OSPF的邻居？什么是OSPF的邻接关系？
    \item OSPF报文的5种类型是什么？每一种类型的用途是什么？
    \item 什么是LSA？怎样区分一个LSA和一个OSPF更新报文？
    \item 简述Hello报文协议的作用
    \item 在多区域OSPF实验中，如果在4个路由器上同时也配置了RIPv2协议，路由表是否会有变化？为什么？
\end{enumerate}

%========================================
\section{实验考试题型总结}
%========================================

\subsection{选择题常考点}

\begin{keybox}
\textbf{网络设备相关}：
\begin{itemize}[itemsep=0pt]
    \item 集线器工作在物理层，交换机工作在数据链路层，路由器工作在网络层
    \item 集线器所有端口共享冲突域，交换机每端口独立冲突域
    \item 只有路由器能隔离广播域
    \item 交换机根据MAC地址转发，路由器根据IP地址转发
\end{itemize}

\textbf{VLAN相关}：
\begin{itemize}[itemsep=0pt]
    \item VLAN可以减小广播域、提高安全性
    \item 不同VLAN之间通信需要三层设备
    \item Trunk端口用于交换机互连，允许多个VLAN通过
    \item Access端口用于连接终端设备
\end{itemize}

\textbf{路由相关}：
\begin{itemize}[itemsep=0pt]
    \item 静态路由需要手动配置，动态路由自动学习
    \item RIP使用跳数作为度量，最大15跳
    \item OSPF使用开销作为度量，收敛速度快
    \item 管理距离：直连0 > 静态1 > OSPF 110 > RIP 120
\end{itemize}
\end{keybox}

\subsection{填空题常考点}

\begin{keybox}
\textbf{协议端口}：
\begin{itemize}[itemsep=0pt]
    \item DNS：UDP 53
    \item FTP：TCP 20(数据)/21(控制)
    \item HTTP：TCP 80
    \item HTTPS：TCP 443
    \item Telnet：TCP 23
    \item SMTP：TCP 25
    \item DHCP：UDP 67(服务器)/68(客户端)
\end{itemize}

\textbf{子网掩码}：
\begin{itemize}[itemsep=0pt]
    \item /24 = 255.255.255.0
    \item /16 = 255.255.0.0
    \item /8 = 255.0.0.0
    \item /22 = 255.255.252.0
\end{itemize}

\textbf{双绞线}：
\begin{itemize}[itemsep=0pt]
    \item 10/100Mbps使用1、2、3、6号线
    \item 最大传输距离100米
    \item T568A：白绿、绿、白橙、蓝、白蓝、橙、白棕、棕
    \item T568B：白橙、橙、白绿、蓝、白蓝、绿、白棕、棕
\end{itemize}
\end{keybox}

\subsection{综合题常考点}

\begin{keybox}
\textbf{IP地址计算}：
\begin{itemize}[itemsep=0pt]
    \item 网络地址 = IP地址 AND 子网掩码
    \item 反掩码 = 255.255.255.255 - 子网掩码
    \item 广播地址 = 网络地址 OR 反掩码
\end{itemize}

\textbf{NAT转换}：
\begin{itemize}[itemsep=0pt]
    \item 出站：源IP从内部本地转换为内部全局
    \item 入站：目的IP从内部全局转换为内部本地
    \item PAT(NAPT)：多个内网IP共享一个公网IP，通过端口区分
\end{itemize}

\textbf{TCP三次握手}：
\begin{itemize}[itemsep=0pt]
    \item 第一次：SYN=1, ACK=0, Seq=x
    \item 第二次：SYN=1, ACK=1, Seq=y, Ack=x+1
    \item 第三次：SYN=0, ACK=1, Seq=x+1, Ack=y+1
\end{itemize}
\end{keybox}

\subsection{实验操作题常考点}

\begin{keybox}
\textbf{VLAN配置}：
\begin{lstlisting}
vlan 10                           ! 创建VLAN
interface f0/1                    ! 进入接口
switchport access vlan 10         ! 划入VLAN
switchport mode trunk             ! 配置Trunk
\end{lstlisting}

\textbf{路由器接口配置}：
\begin{lstlisting}
interface f0/0                    ! 进入接口
ip address 192.168.1.1 255.255.255.0  ! 配置IP
no shutdown                       ! 启用接口
\end{lstlisting}

\textbf{静态路由配置}：
\begin{lstlisting}
ip route 目的网络 子网掩码 下一跳地址
ip route 0.0.0.0 0.0.0.0 下一跳地址  ! 默认路由
\end{lstlisting}

\textbf{ACL配置}：
\begin{lstlisting}
access-list 1 permit 192.168.1.0 0.0.0.255  ! 标准ACL
access-list 101 deny tcp any any eq 80      ! 扩展ACL
ip access-group 1 in                        ! 应用ACL
\end{lstlisting}

\textbf{NAT配置}：
\begin{lstlisting}
ip nat inside                     ! 内部接口
ip nat outside                    ! 外部接口
ip nat pool 名称 起始IP 结束IP netmask 掩码
ip nat inside source list ACL号 pool 名称 overload
\end{lstlisting}
\end{keybox}

%========================================
\section{常见错误与注意事项}
%========================================

\begin{exambox}[配置常见错误]
\textbf{1. 接口未启用}：
\begin{itemize}[itemsep=0pt]
    \item 配置接口后忘记执行\texttt{no shutdown}
    \item 接口状态显示administratively down
\end{itemize}

\textbf{2. 串口时钟未配置}：
\begin{itemize}[itemsep=0pt]
    \item DCE端忘记配置\texttt{clock rate}
    \item 导致line protocol is down
\end{itemize}

\textbf{3. 默认网关未配置}：
\begin{itemize}[itemsep=0pt]
    \item 计算机能ping通本网段，但无法访问其他网段
    \item 需要配置默认网关为路由器接口IP
\end{itemize}

\textbf{4. ACL规则顺序错误}：
\begin{itemize}[itemsep=0pt]
    \item ACL按顺序匹配，先匹配的规则先生效
    \item 应该先写具体规则，后写宽泛规则
    \item 末尾有隐含的deny any
\end{itemize}

\textbf{5. NAT接口方向配置错误}：
\begin{itemize}[itemsep=0pt]
    \item inside/outside配置反了
    \item 导致NAT转换不生效
\end{itemize}

\textbf{6. VLAN间通信问题}：
\begin{itemize}[itemsep=0pt]
    \item 忘记配置Trunk端口
    \item 三层交换机忘记执行\texttt{ip routing}
    \item 计算机网关配置错误
\end{itemize}
\end{exambox}

%========================================
\section{快速查询命令表}
%========================================

\subsection{交换机命令}

\begin{table}[htbp]
\centering\small
\begin{tabular}{ll}
\toprule
\textbf{功能} & \textbf{命令} \\
\midrule
创建VLAN & \texttt{vlan 编号} \\
命名VLAN & \texttt{name 名称} \\
端口划入VLAN & \texttt{switchport access vlan 编号} \\
配置Trunk & \texttt{switchport mode trunk} \\
查看VLAN & \texttt{show vlan} \\
启用生成树 & \texttt{spanning-tree} \\
设置优先级 & \texttt{spanning-tree priority 值} \\
创建聚合端口 & \texttt{interface aggregatePort 编号} \\
加入聚合组 & \texttt{port-group 编号} \\
查看MAC表 & \texttt{show mac-address-table} \\
\bottomrule
\end{tabular}
\caption{交换机常用命令}
\end{table}

\subsection{路由器命令}

\begin{table}[htbp]
\centering\small
\begin{tabular}{ll}
\toprule
\textbf{功能} & \textbf{命令} \\
\midrule
进入接口 & \texttt{interface 类型 编号} \\
配置IP & \texttt{ip address IP 掩码} \\
启用接口 & \texttt{no shutdown} \\
配置时钟 & \texttt{clock rate 速率} \\
静态路由 & \texttt{ip route 网络 掩码 下一跳} \\
默认路由 & \texttt{ip route 0.0.0.0 0.0.0.0 下一跳} \\
查看路由表 & \texttt{show ip route} \\
查看接口 & \texttt{show ip interface brief} \\
查看配置 & \texttt{show running-config} \\
保存配置 & \texttt{copy running-config startup-config} \\
\bottomrule
\end{tabular}
\caption{路由器常用命令}
\end{table}

\subsection{ACL与NAT命令}

\begin{table}[htbp]
\centering\small
\begin{tabular}{ll}
\toprule
\textbf{功能} & \textbf{命令} \\
\midrule
标准ACL & \texttt{access-list 1-99 permit/deny 源地址 反掩码} \\
扩展ACL & \texttt{access-list 100-199 permit/deny 协议 源 目的} \\
应用ACL & \texttt{ip access-group 编号 in/out} \\
查看ACL & \texttt{show access-list} \\
NAT内部接口 & \texttt{ip nat inside} \\
NAT外部接口 & \texttt{ip nat outside} \\
NAT地址池 & \texttt{ip nat pool 名称 起始IP 结束IP netmask 掩码} \\
NAT转换 & \texttt{ip nat inside source list ACL pool 名称 overload} \\
查看NAT & \texttt{show ip nat translations} \\
\bottomrule
\end{tabular}
\caption{ACL与NAT常用命令}
\end{table}

\subsection{Windows网络命令}

\begin{table}[htbp]
\centering\small
\begin{tabular}{ll}
\toprule
\textbf{功能} & \textbf{命令} \\
\midrule
查看IP配置 & \texttt{ipconfig /all} \\
测试连通性 & \texttt{ping 目标IP} \\
路由跟踪 & \texttt{tracert 目标IP} \\
DNS解析 & \texttt{nslookup 域名} \\
清空DNS缓存 & \texttt{ipconfig /flushdns} \\
查看ARP缓存 & \texttt{arp -a} \\
清空ARP缓存 & \texttt{arp -d} \\
查看路由表 & \texttt{route print} \\
查看网络连接 & \texttt{netstat -an} \\
\bottomrule
\end{tabular}
\caption{Windows网络命令}
\end{table}

\subsection{RIP与OSPF命令}

\begin{table}[htbp]
\centering\small
\begin{tabular}{ll}
\toprule
\textbf{功能} & \textbf{命令} \\
\midrule
启用RIP & \texttt{router rip} \\
RIP版本2 & \texttt{version 2} \\
RIP网络 & \texttt{network 网络地址} \\
启用OSPF & \texttt{router ospf} \\
OSPF网络 & \texttt{network 网络 反掩码 area 区域号} \\
查看路由协议 & \texttt{show ip protocols} \\
查看OSPF邻居 & \texttt{show ip ospf neighbor} \\
查看OSPF数据库 & \texttt{show ip ospf database} \\
\bottomrule
\end{tabular}
\caption{RIP与OSPF命令}
\end{table}

%========================================
\section{Trunk与Access端口详解}
%========================================

\subsection{端口工作模式}

\begin{keybox}
\textbf{Access端口}：
\begin{itemize}[itemsep=0pt]
    \item 用于连接终端设备（计算机、打印机等）
    \item 只能属于一个VLAN
    \item 发送数据时不带VLAN标签
    \item 接收数据时添加所属VLAN标签
\end{itemize}

\textbf{Trunk端口}：
\begin{itemize}[itemsep=0pt]
    \item 用于交换机之间互连
    \item 可以传输多个VLAN的数据
    \item 发送数据时带VLAN标签（802.1Q）
    \item 接收数据时根据标签转发到对应VLAN
\end{itemize}
\end{keybox}

\subsection{Trunk配置详解}

\begin{cmdbox}
\textbf{配置Trunk端口}：
\begin{lstlisting}
Switch(config)# interface fastEthernet 0/24
Switch(config-if)# switchport mode trunk
! switchport mode trunk: 设置端口为Trunk模式
! Trunk端口允许所有VLAN通过(默认)
Switch(config-if)# exit
\end{lstlisting}

\textbf{重新设置为Access模式}：
\begin{lstlisting}
Switch(config-if)# switchport mode access
! 将端口从Trunk模式改回Access模式
\end{lstlisting}

\textbf{查看VLAN配置}：
\begin{lstlisting}
Switch# show vlan
! 显示VLAN信息和端口分配
! 注意: Trunk端口不会显示在VLAN列表中
\end{lstlisting}
\end{cmdbox}

\begin{keybox}
\textbf{重要说明}：
\begin{itemize}[itemsep=0pt]
    \item 如果没有给VLAN指定名称，系统会自动命名
    \item 查看VLAN配置时，Trunk端口不会显示在任何VLAN的端口列表中
    \item 两台交换机互连时，两端都要配置为Trunk模式
\end{itemize}
\end{keybox}

%========================================
\section{广播风暴与环路问题}
%========================================

\subsection{环路产生原因}

\begin{keybox}
当交换机之间有多条链路连接时，会形成环路，造成以下问题：
\begin{itemize}[itemsep=0pt]
    \item \textbf{广播风暴}：广播帧在环路中无限循环
    \item \textbf{MAC地址表不稳定}：同一MAC地址从不同端口学习
    \item \textbf{重复帧}：同一帧被多次接收
    \item \textbf{网络瘫痪}：交换机CPU占用过高，无法正常工作
\end{itemize}

\textbf{现象}：
\begin{itemize}[itemsep=0pt]
    \item 交换机配置无法操作，类似死机现象
    \item 交换机端口指示灯频繁闪烁
    \item 断开任一条连线，恢复正常
\end{itemize}
\end{keybox}

\subsection{解决方案}

\begin{itemize}[leftmargin=2em,itemsep=0pt]
    \item \textbf{生成树协议(STP)}：自动阻塞冗余链路，消除环路
    \item \textbf{链路聚合}：将多条物理链路捆绑为一条逻辑链路
\end{itemize}

\begin{keybox}
\textbf{注意}：在生成树或链路聚合配置之前，不要把交换机间两条连线全连上！
\end{keybox}

%========================================
\section{生成树协议深入理解}
%========================================

\subsection{STP工作原理}

\begin{keybox}
\textbf{生成树协议选举过程}：

\textbf{1. 选举根桥(Root Bridge)}：
\begin{itemize}[itemsep=0pt]
    \item 比较Bridge ID（优先级 + MAC地址）
    \item 优先级越小越优先（默认32768）
    \item 优先级相同则比较MAC地址，越小越优先
\end{itemize}

\textbf{2. 选举根端口(Root Port)}：
\begin{itemize}[itemsep=0pt]
    \item 非根桥上到根桥开销最小的端口
    \item 开销相同则选端口优先级最小的
    \item 每个非根桥只有一个根端口
\end{itemize}

\textbf{3. 选举指定端口(Designated Port)}：
\begin{itemize}[itemsep=0pt]
    \item 每个网段上到根桥开销最小的端口
    \item 根桥的所有端口都是指定端口
\end{itemize}

\textbf{4. 阻塞非指定端口}：
\begin{itemize}[itemsep=0pt]
    \item 既不是根端口也不是指定端口的端口被阻塞
    \item 阻塞端口不转发数据，但接收BPDU
\end{itemize}
\end{keybox}

\subsection{STP端口状态}

\begin{table}[htbp]
\centering
\begin{tabular}{lcc}
\toprule
\textbf{状态} & \textbf{学习MAC} & \textbf{转发数据} \\
\midrule
Blocking（阻塞） & 否 & 否 \\
Listening（监听） & 否 & 否 \\
Learning（学习） & 是 & 否 \\
Forwarding（转发） & 是 & 是 \\
Disabled（禁用） & 否 & 否 \\
\bottomrule
\end{tabular}
\caption{STP端口状态}
\end{table}

\subsection{show spanning-tree输出解读}

\begin{cmdbox}
\begin{lstlisting}
Switch# show spanning-tree
StpVersion: RSTP              ! 生成树协议版本
SysStpStatus: Enabled         ! 生成树状态：已启用
BaseNumPorts: 24              ! 端口数量
MaxAge: 20                    ! 最大老化时间(秒)
HelloTime: 2                  ! Hello报文间隔(秒)
ForwardDelay: 15              ! 转发延迟(秒)
BridgeAddr: 00d0.f87b.ed31    ! 交换机MAC地址
Priority: 0                   ! 优先级(已设置为最高)
DesignatedRoot: 000000D0F87BED31  ! 根桥ID
RootCost: 0                   ! 到根桥的开销(0表示自己是根桥)
RootPort: 0                   ! 根端口(0表示自己是根桥)
\end{lstlisting}

\textbf{判断根桥}：
\begin{itemize}[itemsep=0pt]
    \item RootCost = 0 表示自己是根桥
    \item RootPort = 0 表示自己是根桥
    \item DesignatedRoot = 自己的BridgeAddr
\end{itemize}
\end{cmdbox}

%========================================
\section{链路聚合深入理解}
%========================================

\subsection{链路聚合优势}

\begin{keybox}
\textbf{链路聚合(Link Aggregation)的优势}：
\begin{itemize}[itemsep=0pt]
    \item \textbf{增加带宽}：多条链路带宽叠加
    \item \textbf{提高可靠性}：一条链路故障，其他链路继续工作
    \item \textbf{负载均衡}：流量分散到多条链路
    \item \textbf{无需STP阻塞}：聚合后视为一条逻辑链路，不会形成环路
\end{itemize}
\end{keybox}

\subsection{流量均衡算法}

\begin{keybox}
\textbf{流量均衡方式}：
\begin{itemize}[itemsep=0pt]
    \item \textbf{src-mac}：根据源MAC地址分配（默认）
    \item \textbf{dst-mac}：根据目的MAC地址分配
    \item \textbf{src-ip}：根据源IP地址分配
    \item \textbf{dst-ip}：根据目的IP地址分配
\end{itemize}

\textbf{配置命令}：
\begin{lstlisting}
Switch# aggregateport load-balance dst-mac
! 按目的MAC地址进行流量分配
Switch# aggregateport load-balance src-ip
! 按源IP地址进行流量分配
\end{lstlisting}
\end{keybox}

\subsection{链路聚合与生成树对比}

\begin{table}[htbp]
\centering
\begin{tabular}{lcc}
\toprule
\textbf{特性} & \textbf{生成树} & \textbf{链路聚合} \\
\midrule
带宽利用 & 只用一条链路 & 所有链路同时使用 \\
故障切换 & 有短暂中断 & 几乎无中断 \\
配置复杂度 & 简单 & 稍复杂 \\
适用场景 & 冗余备份 & 高带宽需求 \\
\bottomrule
\end{tabular}
\caption{生成树与链路聚合对比}
\end{table}

%========================================
\section{静态路由配置详解}
%========================================

\subsection{静态路由命令格式}

\begin{cmdbox}
\textbf{静态路由配置格式}：
\begin{lstlisting}
Router(config)# ip route 网络地址 子网掩码 下一跳地址
! 或
Router(config)# ip route 网络地址 子网掩码 出接口
\end{lstlisting}

\textbf{参数说明}：
\begin{itemize}[itemsep=0pt]
    \item 网络地址：目的网络的网络号
    \item 子网掩码：目的网络的子网掩码
    \item 下一跳地址：转发数据包的下一个路由器接口IP
    \item 出接口：本地路由器转发数据的接口
\end{itemize}

\textbf{示例}：
\begin{lstlisting}
! 方式1：指定下一跳IP地址
RA(config)# ip route 17.3.0.0 255.255.0.0 17.2.0.2
! 到达17.3.0.0/16网络，下一跳为17.2.0.2

! 方式2：指定出接口
RA(config)# ip route 17.3.0.0 255.255.0.0 serial 1/3
! 到达17.3.0.0/16网络，从serial 1/3接口转发
\end{lstlisting}
\end{cmdbox}

\subsection{默认路由}

\begin{keybox}
\textbf{默认路由}：当路由表中没有匹配的路由时使用。

\begin{lstlisting}
Router(config)# ip route 0.0.0.0 0.0.0.0 下一跳地址
! 0.0.0.0 0.0.0.0 表示匹配所有网络
\end{lstlisting}

\textbf{应用场景}：
\begin{itemize}[itemsep=0pt]
    \item 连接Internet的边界路由器
    \item 只有一个出口的网络
    \item 简化路由配置
\end{itemize}
\end{keybox}

%========================================
\section{ACL应用位置原则}
%========================================

\begin{keybox}
\textbf{ACL应用位置原则}：

\textbf{标准ACL}：
\begin{itemize}[itemsep=0pt]
    \item 只能根据源地址过滤
    \item 应用在\textbf{靠近目的地}的位置
    \item 原因：避免过早丢弃数据包，影响其他目的地的访问
\end{itemize}

\textbf{扩展ACL}：
\begin{itemize}[itemsep=0pt]
    \item 可以根据源地址、目的地址、协议、端口过滤
    \item 应用在\textbf{靠近源地址}的位置
    \item 原因：尽早丢弃不需要的数据包，节省网络带宽
\end{itemize}
\end{keybox}

\subsection{ACL方向说明}

\begin{keybox}
\textbf{in和out的含义}：
\begin{itemize}[itemsep=0pt]
    \item \textbf{in}：检查进入该接口的数据流量
    \item \textbf{out}：检查离开该接口的数据流量
\end{itemize}

\textbf{示例}：
\begin{lstlisting}
! 在F1/0接口上，检查出方向的流量
Router(config)# interface fastEthernet 1/0
Router(config-if)# ip access-group 101 out
! 数据包从路由器通过F1/0发出时被检查
\end{lstlisting}
\end{keybox}

%========================================
\section{NAT地址术语}
%========================================

\begin{keybox}
\textbf{NAT四种地址类型}：

\begin{center}
\begin{tabular}{lll}
\toprule
\textbf{术语} & \textbf{含义} & \textbf{示例} \\
\midrule
内部本地地址 & 内网主机的私有IP & 192.168.1.3 \\
内部全局地址 & 转换后的公网IP & 59.67.148.3 \\
外部本地地址 & 外网主机在内网看到的IP & 202.113.64.2 \\
外部全局地址 & 外网主机的真实IP & 202.113.64.2 \\
\bottomrule
\end{tabular}
\end{center}

\textbf{NAT转换过程}：
\begin{itemize}[itemsep=0pt]
    \item 出站：源IP从内部本地 → 内部全局
    \item 入站：目的IP从内部全局 → 内部本地
\end{itemize}
\end{keybox}

%========================================
\section{路由器接口类型详解}
%========================================

\subsection{常见接口类型}

\begin{table}[htbp]
\centering
\begin{tabular}{llll}
\toprule
\textbf{接口类型} & \textbf{缩写} & \textbf{速率} & \textbf{用途} \\
\midrule
FastEthernet & F/Fa & 100Mbps & 局域网连接 \\
GigabitEthernet & G/Gi & 1Gbps & 高速局域网 \\
Serial & S & 可配置 & 广域网连接 \\
Console & Con & - & 本地管理 \\
\bottomrule
\end{tabular}
\caption{路由器接口类型}
\end{table}

\subsection{接口编号格式}

\begin{keybox}
\textbf{接口编号格式}：类型 插槽号/端口号

\textbf{示例}：
\begin{itemize}[itemsep=0pt]
    \item \texttt{fastEthernet 1/0}：插槽1的第0个快速以太网口
    \item \texttt{serial 1/3}：插槽1的第3个串口
    \item \texttt{gigabitEthernet 0/0}：插槽0的第0个千兆以太网口
\end{itemize}
\end{keybox}

%========================================
\section{实验拓扑IP规划方法}
%========================================

\begin{keybox}
\textbf{IP地址规划原则}：

\textbf{1. 网段划分}：
\begin{itemize}[itemsep=0pt]
    \item 每个LAN网段使用独立的子网
    \item 路由器之间的串口连接使用独立的子网
    \item 子网掩码根据主机数量选择
\end{itemize}

\textbf{2. 地址分配}：
\begin{itemize}[itemsep=0pt]
    \item 路由器接口通常使用网段的第一个或最后一个可用地址
    \item 计算机使用网段中间的地址
    \item 预留部分地址用于扩展
\end{itemize}

\textbf{3. 网关配置}：
\begin{itemize}[itemsep=0pt]
    \item 计算机的默认网关 = 所在网段的路由器接口IP
    \item 不同网段的计算机网关不同
\end{itemize}
\end{keybox}

\subsection{IP规划示例}

\begin{exambox}[典型实验IP规划]
\textbf{场景}：两台路由器RA、RB，两台PC

\textbf{IP规划}（使用202.115.0.0/16）：
\begin{itemize}[itemsep=0pt]
    \item PC1网段：202.115.10.0/24
    \item PC2网段：202.115.20.0/24
    \item 串口互连：202.115.30.0/24
\end{itemize}

\textbf{地址分配}：
\begin{itemize}[itemsep=0pt]
    \item RA的F1/0：202.115.10.1（PC1的网关）
    \item RA的S1/3：202.115.30.1
    \item RB的S1/3：202.115.30.2
    \item RB的F1/0：202.115.20.1（PC2的网关）
    \item PC1：202.115.10.2，网关202.115.10.1
    \item PC2：202.115.20.2，网关202.115.20.1
\end{itemize}
\end{exambox}

%========================================
\section{Wireshark抓包分析要点}
%========================================

\subsection{协议层次识别}

\begin{keybox}
\textbf{Wireshark显示的协议层次}：

\begin{itemize}[itemsep=0pt]
    \item \textbf{Frame}：物理层帧信息
    \item \textbf{Ethernet II}：数据链路层（源/目的MAC）
    \item \textbf{Internet Protocol}：网络层（源/目的IP）
    \item \textbf{TCP/UDP}：传输层（源/目的端口）
    \item \textbf{HTTP/DNS/FTP等}：应用层协议
\end{itemize}
\end{keybox}

\subsection{常见抓包场景分析}

\begin{exambox}[Ping命令抓包]
\textbf{执行}：\texttt{ping 192.168.1.1}

\textbf{抓包特征}：
\begin{itemize}[itemsep=0pt]
    \item 协议：ICMP
    \item 类型：Echo Request（请求）/ Echo Reply（回复）
    \item 无传输层（ICMP直接封装在IP中）
\end{itemize}
\end{exambox}

\begin{exambox}[FTP连接抓包]
\textbf{执行}：FTP连接服务器

\textbf{抓包特征}：
\begin{itemize}[itemsep=0pt]
    \item 协议：TCP + FTP
    \item 控制连接：目的端口21
    \item 数据连接：目的端口20
    \item 可以看到用户名和密码（明文传输）
\end{itemize}
\end{exambox}

\begin{exambox}[ARP请求抓包]
\textbf{场景}：首次访问同网段主机

\textbf{抓包特征}：
\begin{itemize}[itemsep=0pt]
    \item 协议：ARP
    \item 目的MAC：ff:ff:ff:ff:ff:ff（广播）
    \item Opcode：1（请求）/ 2（回复）
    \item 包含目标IP和发送方IP/MAC
\end{itemize}
\end{exambox}

%========================================
\section{实验操作注意事项}
%========================================

\begin{keybox}
\textbf{实验前准备}：
\begin{enumerate}[itemsep=0pt]
    \item 检查设备连线是否正确
    \item 确认Console线连接正常
    \item 打开终端软件（如SecureCRT）
    \item 设置正确的串口参数（9600, 8, N, 1）
\end{enumerate}

\textbf{配置过程注意}：
\begin{enumerate}[itemsep=0pt]
    \item 每配置一步，及时验证
    \item 使用show命令查看配置结果
    \item 配置完成后保存配置
    \item 记录配置命令，便于排错
\end{enumerate}

\textbf{故障排查步骤}：
\begin{enumerate}[itemsep=0pt]
    \item 检查物理连接（网线、接口指示灯）
    \item 检查接口状态（show ip interface brief）
    \item 检查IP配置（ipconfig、show running-config）
    \item 检查路由配置（show ip route）
    \item 使用ping逐段测试连通性
\end{enumerate}
\end{keybox}

%========================================
\section{私有IP地址范围}
%========================================

\begin{keybox}
\textbf{RFC 1918定义的私有IP地址}：

\begin{center}
\begin{tabular}{lll}
\toprule
\textbf{类别} & \textbf{地址范围} & \textbf{CIDR表示} \\
\midrule
A类 & 10.0.0.0 - 10.255.255.255 & 10.0.0.0/8 \\
B类 & 172.16.0.0 - 172.31.255.255 & 172.16.0.0/12 \\
C类 & 192.168.0.0 - 192.168.255.255 & 192.168.0.0/16 \\
\bottomrule
\end{tabular}
\end{center}

\textbf{特点}：
\begin{itemize}[itemsep=0pt]
    \item 私有地址不能在Internet上路由
    \item 内网使用私有地址，通过NAT访问外网
    \item 不同组织可以使用相同的私有地址
\end{itemize}
\end{keybox}

%========================================
\section{子网划分计算}
%========================================

\subsection{子网划分公式}

\begin{keybox}
\textbf{子网划分计算}：

\begin{itemize}[itemsep=0pt]
    \item 子网数 = $2^n$（n为借用的主机位数）
    \item 每个子网主机数 = $2^m - 2$（m为剩余主机位数，减2是去掉网络地址和广播地址）
\end{itemize}

\textbf{常用子网掩码}：
\begin{itemize}[itemsep=0pt]
    \item /24 = 255.255.255.0，主机数254
    \item /25 = 255.255.255.128，主机数126
    \item /26 = 255.255.255.192，主机数62
    \item /27 = 255.255.255.224，主机数30
    \item /28 = 255.255.255.240，主机数14
    \item /30 = 255.255.255.252，主机数2（点对点链路常用）
\end{itemize}
\end{keybox}

\subsection{子网计算示例}

\begin{exambox}[子网划分题]
\textbf{题目}：将192.168.1.0/24划分为4个子网

\textbf{解答}：
\begin{itemize}[itemsep=0pt]
    \item 需要借用2位（$2^2=4$）
    \item 新子网掩码：/26 = 255.255.255.192
    \item 子网1：192.168.1.0/26（可用：.1-.62）
    \item 子网2：192.168.1.64/26（可用：.65-.126）
    \item 子网3：192.168.1.128/26（可用：.129-.190）
    \item 子网4：192.168.1.192/26（可用：.193-.254）
\end{itemize}
\end{exambox}

\end{document}
